\chapter{Modèles de ZFC}
\label{chp.modZFC}

\minitoc

\lettrine{N}{ous} avons vu dans les chapitres précédents combien $\ZFC$ était
une théorie riche en terme d'expressivité, au point même où l'ensemble des
mathématiques peut se formaliser dans $\ZFC$. Il devient alors assez technique
de parler de modèle de $\ZFC$, puisqu'un tel modèle est en quelque sorte un
modèle de toutes les mathématiques que l'on connaît.

Une première conséquence de cette expressivité est un certain flou entre ce que
l'on exprime dans $\ZFC$ et dans les mathématiques usuelles. Puisque tout peut
se définir en terme d'ensembles, il semble que tout ce dont nous avons parlé
jusque là puisse se formaliser dans $\ZFC$. C'est le cas, mais il faut faire
attention à ce qu'on manipule.

En particulier, un modèle de $\ZFC$, nommons-le $\varMM$, contient lui-même
un ensemble $\Formula(\Sign)$ pour une signature $\Sign$ donnée. On pourrait
alors se dire que, comme on peut formaliser nos résultats de logique en terme
d'ensembles, on peut travailler au sein de $\ZFC$ pour faire de la logique.
Malheureusement, les propositions de $\Formula(\Sign)$ ne sont pas les mêmes
que celles que nous connaissons~: on peut facilement interpréter une proposition
habituelle par un élément de $\Formula(\Sign)$, mais l'inverse n'a pas de
raison d'être vrai.

Cela mène à deux résultats essentiels pour décrire les limitations de cette
approche formelle de la logique~: les théorèmes d'incomplétudes de Gödel
(\cref{thm.Godel1.ZFC} et \cref{thm.Godel2.ZFC}) et le théorème
d'indéfinissabilité de la vérité de Tarski (\cref{thm.Tarski.ZFC}), que nous
abordons dans le \cref{chp.arith}. Ceux-ci ont pour conséquence que si l'on
considère un codage $\godcod{\tok}$ qui à une formule au sens habituel associe
une formule de notre modèle $\varMM$, alors il n'existe pas de formule
$\varFF_{\varMM}$ à l'intérieur de $\varMM$ telle que
\[\forall \varFP \in \Formula(\Sign),
\varMM \satisf \varFF_{\varMM}(\godcod{\varFP}) \iff \varFP\]

Cette limitation implique que l'on ne peut pas à strictement
parler écrire une formule comme $\unZF\satisf \ZFC$ pour notre univers
ensembliste $\unZF$. Par contre, il est possible d'écrire par exemple
$\ZFC\vdP \varFF$, qui en utilisant le théorème de complétude permet de dire
que $\unZF\satisf \varFF$~: ainsi, l'écriture de $\unZF\satisf\ZFC$
se traduit en fait par une infinité de théorèmes de la forme
$\ZFC\vD\varFF$ pour chaque axiome $\varFF \in \ZFC$.

S'il est assez clair que $\ZFC\vD\ZFC$ est un constat assez vide, cela nous
permet ensuite d'introduire la relativisation, qui permet d'exprimer la notion
de modèle classe de $\ZFC$.

Ce chapitre s'attachera à l'étude des outils basiques à propos des modèles
classe. Nous verrons d'abord les définitions et conséquences de cette notion,
puis nous aborderons l'absoluité, qui permet d'étudier la conservation entre des
modèles de $\ZFC$.

Nous verrons ensuite la théorie des modèles intérieurs, et son exemple le plus
connu~: l'univers constructible de Gödel $\bL$.

\section{Premiers pas}

\subsection{Relativisation}

On commence par introduire la notion de relativisation d'une formule. Comme on
l'a dit, il est impossible de représenter de façon interne à un modèle de $\ZFC$
que celui-ci vérifie une formule~: il faut donc parler de façon externe au
modèle. Pour exprimer que $\varFF$ est vraie dans l'univers ambiant, il suffit
alors de travailler normalement, en utilisant les axiomes de $\ZFC$. Cependant,
cette approche est très limitée, et on aimerait s'autoriser à parler d'autres
classes que l'univers $\unZF$ tout entier. Pour ce faire, comme on doit
procéder de façon purement syntaxique, on modifie la formule même qu'on veut
prouver en y incluant l'information de la classe qu'on utilise.

Disons qu'on veuille montrer qu'une formule $\varFF$ est vraie dans la classe
des ordinaux, l'idée est simplement de modifier tous les quantificateurs pour
que ceux-ci ne parlent que d'éléments de $\Ord$. La relativisation est
précisément ce procédé de réduire la portée des quantificateurs par une formule,
exprimant en général l'appartenance à une classe qui nous intéresse.

\begin{definition}[Relativisation d'une formule]\label{def.relativisation}
  Soient $\varFF,\varFP\in\Formula(\SignP{\ZF})$ où $\varFP$ possède une
  variable libre. On appelle relativisation de $\varFF$ à $\varFP$ la formule
  $\varFF^\varFP$ définie par induction sur $\varFF$ par~:
  \begin{itemize}
  \item si $\varFF$ est atomique alors $\varFF^\varFP \defeq \varFF$.
  \item si $\varFF = \top$ alors $\varFF^\varFP \defeq \top$.
  \item si $\varFF = \bot$ alors $\varFF^\varFP \defeq \bot$.
  \item si $\varFF = \lnot \varFC$ alors
    $\varFF^\varFP \defeq \lnot \varFC^\varFP$.
  \item si $\varFF = \varFC_1 \lor \varFC_2$ alors
    $\varFF^\varFP \defeq {\varFC_1}^\varFP \lor {\varFC_2}^\varFP$.
  \item si $\varFF = \varFC_1 \land \varFC_2$ alors
    $\varFF^\varFP \defeq {\varFC_1}^\varFP \land {\varFC_2}^\varFP$.
  \item si $\varFF = \varFC_1 \to \varFC_2$ alors
    $\varFF^\varFP \defeq {\varFC_1}^\varFP \to {\varFC_2}^\varFP$.
  \item si $\varFF = \exists x, \varFC$ alors
    $\varFF^\varFP \defeq \exists x, \varFP(x) \land \varFC^\varFP$.
  \item si $\varFF = \forall x, \varFC$ alors
    $\varFF^\varFP \defeq \forall x, \varFP(x) \to \varFC^\varFP$.
  \end{itemize}
\end{definition}

Ce procédé de relativisation se comporte bien vis à vis de la dérivation~: si
$\varCG\vdP\varFF$ est dérivable, alors $\varCG^\varFP\vdP\varFF^\varFP$ aussi.
Cependant, il faut faire attention ici à la gestion des variables libres~:
pour que l'induction fonctionne correctement, on doit ajouter au contexte
$\varCG$ les hypothèses $\varFP(x_1),\ldots,\varFP(x_n)$ pour $x_1,\ldots,x_n$
les variables libres dans le séquent.

\begin{lemma}\label{lem.relat.deduc}
  Pour $\varCG\in\Ctx(\SignP{\ZF})$ et
  $\varFF,\varFP\in\Formula(\SignP{\ZF})$, si $\varCG\vdP \varFF$ alors, en
  notant $\{x_1,\ldots,x_n\}$ l'ensemble $\VL(\varCG)\cup\VL(\varFF)$, et
  $\varFP(\vec x)$ le contexte $\varFP(x_1),\ldots,\varFP(x_n)$, on a
  $\varFP(\vec x),\varCG^\varFP\vdP\varFF^\varFP$.
\end{lemma}

\begin{proof}
  On procède par induction sur la relation $\varCG \vdP \varFF$~:
  \begin{itemize}
  \item[$\Rax$~:] si $\varFF \in \varCG$, alors $\varFF^\varFP \in \varCG^\varFP$,
    d'où le résultat par affaiblissement.
  \item[$\RtopI$~:] peu importe $\varCG$, on a toujours
    $\varFP(\vec x),\varCG^\varFP \vdP \top$.
  \item[$\RbotC$~:] on suppose que
    $\varFP(\vec x),\varCG^\varFP, \lnot \varFF^\varFP \vdP \bot$, alors on peut
    appliquer $\RbotC$ pour en déduire que
    $\varFP(\vec x), \varCG^\varFP \vdP \varFF^\varFP$.
  \item[$\RorI$~:] on suppose que
    $\varFP(\vec x),\varCG^\varFP \vdP {\varFF_i}^\varFP$ pour $i \in \{1,2\}$,
    donc
    $\varFP(\vec x),\varCG^\varFP \vdP {\varFF_1}^\varFP\lor {\varFF_2}^\varFP$,
    c'est-à-dire
    $\varFP(\vec x),\varCG^\varFP \vdP (\varFF_1\lor \varFF_2)^\varFP$.
  \item[$\RorE$~:] on suppose que
    $\varFP(\vec x),\varCG^\varFP\vdP {\varFF_1}^\varFP\lor {\varFF_2}^\varFP$,
    que $\varFP(\vec x),\varCG^\varFP, {\varFF_1}^\varFP \vdP \varFC^\varFP$
    et que $\varFP(\vec x),\varCG^\varFP,{\varFF_2}^\varFP\vdP \varFC^\varFP$,
    alors on peut appliquer la règle
    $\RorE$ pour en déduire que
    $\varFP(\vec x),\varCG^\varFP \vdP \varFC^\varFP$.
  \item[$\RandI$~:] on suppose que
    $\varFP(\vec x),\varCG^\varFP\vdP {\varFF_1}^\varFP$ et
    $\varFP(\vec x),\varCG^\varFP\vdP {\varFF_2}^\varFP$, alors
    $\varFP(\vec x),\varCG^\varFP\vdP {\varFF_1}^\varFP \land {\varFF_2}^\varFP$.
  \item[$\RandE$~:] on suppose que
    $\varFP(\vec x),\varCG^\varFP\vdP (\varFF_1\land \varFF_2)^\varFP$,
    c'est-à-dire que
    $\varFP(\vec x),\varCG^\varFP\vdP {\varFF_1}^\varFP \land {\varFF_2}^\varFP$,
    alors on peut en déduire que
    $\varFP(\vec x),\varCG^\varFP\vdP {\varFF_i}^\varFP$ pour
    $i \in \{1,2\}$.
  \item[$\RtoI$~:] si
    $\varFP(\vec x),\varCG^\varFP,{\varFF_1}^\varFP\vdP \varFF^\varFP$ alors
    $\varFP(\vec x),\varCG^\varFP\vdP {\varFF_1}^\varFP \to {\varFF_2}^\varFP$, et
    par définition on sait que
    ${\varFF_1}^\varFP\to {\varFF_2}^\varFP = (\varFF_1\to \varFF_2)^\varFP$, d'où
    le résultat.
  \item[$\RtoE$~:] si
    $\varFP(\vec x),\varCG^\varFP \vdP (\varFF_1\to\varFF_2)^\varFP$ et
    $\varFP(\vec x),\varCG^\varFP \vdP {\varFF_1}^\varFP$, alors
    $\varFP(\vec x),\varCG^\varFP\vdP {\varFF_2}^\varFP$.
  \item[$\RforaI$~:] si
    $\varFP(\vec x),\varCG^\varFP\vdP \varFF[v/x]^\varFP$, où $v$ est
    libre dans $\varCG$ et $\varFF$, alors par affaiblissement on sait que
    $\varFP(v),\varFP(\vec x),\varCG^\varFP\vdP\varFF[v/x]^\varFP$, donc que
    $\varFP(\vec x),\varCG^\varFP\vdP\varFP(v)\to \varFF[v/x]^\varFP$, d'où par la
    règle
    $\RforaI$ le fait que
    $\varFP(\vec x),\varCG^\varFP \vdP (\forall x,\varFF)^\varFP$.
  \item[$\RforaE$~:] si
    $\varFP(\vec x),\varCG^\varFP \vdP (\forall x,\varFF)^\varFP$, et
    pour une variable $y \in \VL(\Gamma)\cup \VL(A)$ (puisque tout terme dans
    $\SignP{\ZF}$ est une variable et que la relativisation n'ajoute pas
    de variable libre), on peut spécialiser notre séquent en
    $\varFP(\vec x),\varCG^\varFP\vdP\varFP(y)\to\varFF[y/x]$ et, comme
    $y$ est une variable libre, on sait que
    $\varFP(\vec x),\varCG^\varFP\vdP\varFP(y)$, d'où le résultat.
  \item[$\RexI$~:] si $\varFP(\vec x),\varCG^\varFP \vdP \varFF[y/x]$ alors on
    sait que $\varFP(\vec x), \varCG^\varFP\vdP \varFP(y)$, d'où le résultat par
    la règle $\RandI$.
  \item[$\RexE$~:] si
    $\varFP(\vec x),\varCG^\varFP \vdP (\exists x, \varFF)^\varFP$ et
    $\varFP(\vec x),\varCG^\varFP, \varFP(x)\land\varFF^\varFP\vdP\varFC^\varFP$,
    alors la règle $\RexE$ nous donne directement
    $\varFP(\vec x),\varCG^\varFP \vdP \varFC^\varFP$. On
    remarque qu'on peut effectivement utiliser l'hypothèse d'induction, car
    $\varFP(x)\land\varFF^\varFP$ est $\varFP(x),\varFF^\varFP$ sont des
    contextes équivalents.
  \item[$\ReqI$~:] si $t$ est un terme, alors
    $\varFP(\vec x),\varCG^\varFP\vdP t = t$.
  \item[$\ReqE$~:] si $\varFP(\vec x),\varCG^\varFP \vdP t = u$ et
    $\varFP(\vec x), \varCG^\varFP \vdP \varFF[t/x]^\varFP$, alors on peut
    vérifier par induction sur $\varFF$ que
    $(\varFF[t/x])^\varFP = \varFF^\varFP [t/x]$, donc on peut directement
    appliquer $\ReqE$ pour obtenir le résultat.
  \end{itemize}
  Ainsi, par induction, on sait que
  $\varFP(\vec x),\varCG^\varFP \vdP \varFF^\varFP$.
\end{proof}

\begin{remark}
  La preuve précédente serait bien plus naturelle dans le cas de séquents
  avec variables explicites. Sans gestion des variables du premier ordre, ce
  lemme serait faux~: en prenant pour $\varFP$ un prédicat toujours faux, tel
  que $x\neq x$, on peut prouver
  $\vdP (\forall x, \varFF)\to (\exists x, \varFF)$ mais pas
  $\vdP (\forall x,\varFF)^\varFP \to (\exists x, \varFF)^\varFP$ puisque le $x$
  invoqué par la première formule ne vérifie pas $\varFP$. Dans un tel cas, en
  ajoutant la gestion des variables, on a une hypothèse absurde donc la
  conclusion est forcément vraie.
\end{remark}

\begin{remark}
  Il est possible de généraliser ce lemme à une signature $\Sign$ quelconque.
  Cependant, pour cela, il est nécessaire que $\varFP$ vérifie une propriété de
  stabilité~: il faut que pour chaque symbole de fonction $f$ d'arité $n$,
  on puisse prouver
  \[\vdP\forall x_1,\ldots,x_n, \left(\bigwedge_{i = 1}^n \varFP(x_i)\right)
  \to \varFP(f(x_1,\ldots,x_n))\]
  Ce faisant, on peut montrer que tout terme vérifie $\varFP$ (en supposant que
  les variables vérifient $\varFP$), et ainsi vérifier les parties de
  l'induction liées à l'élimination du $\forall$ et l'introduction du $\exists$.
\end{remark}

On peut considérer le cas particulier de la cohérence, qui correspond à la
preuve de $\bot$, pour montrer la propriété suivante.

\begin{theorem}[Préservation de la cohérence]
  Soient $\varT$ et $\varT'$ deux théories sur $\SignP{\ZF}$ et
  $\varFP(x)$ une formule à une variable libre dans ce même langage telles que
  pour toute formule $\varFF \in \varT'$, $\varT\vdP \varFF^\varFP$. Alors
  \[\Coher(\varT) \implies \Coher(\varT')\]
\end{theorem}

\begin{proof}
  On donne deux preuves, l'une syntaxique et l'autre sémantique.
  \begin{itemize}
  \item supposons $\lnot\Coher(\varT')$, cela signifie donc qu'il existe une
    preuve $\varT'\vdP \bot$, et comme cette preuve est finie on peut
    extraire une partie finie $F\subfin \varT'$ telle que $F\vdP \bot$.
    En utilisant le \cref{lem.relat.deduc}, il vient donc que
    $F^\varFP\vdP \bot^\varFP$ (comme les énoncés sont clos, on n'a pas à rajouter
    de $\varFP(x)$ pour des variables $x$). Or, $\bot^\varFP = \bot$, donc on en
    déduit que $F^\varFP\vdP \bot$. Comme $\varT\vdP \varFF^\varFP$ pour
    chaque $\varFF \in F$, on en déduit que $\varT\vdP \bot$, donc que
    $\lnot\Coher(\varT)$. Par contraposée, on en déduit l'implication
    voulue.
  \item supposons $\Coher(\varT)$. On trouve donc un modèle $\varMM$
    de $\varT$. Soit $\varMN$ le sous-modèle de $\varMM$ défini
    par
    \[\carM{\varMN}\defeq \compr{x\in\carMM}{\varMM\satisf\varFP(x)}\]
    Par hypothèse, on sait que $\varMM\satisf{\varT'}^\varFP$, donc
    $\varMN\satisf \varT$, d'où $\Coher(\varT)$.
  \end{itemize}

  Les deux approches nous donnent donc que
  $\Coher(\varT)\implies\Coher(\varT')$.
\end{proof}

On peut par exemple appliquer ce théorème à $\ZFC - \AxF$ et $\ZFC$ pour dire
que s'il existe une classe $\bV$ telle que
$\ZFC - \AxF\vdP \ZFC^{\bV}$, alors les deux théories sont équicohérentes
(puisque dans l'autre sens, il est évident qu'enlever des axiomes ne change pas
la cohérence).

\subsection{L'univers bien fondé de Von Neumann}

On peut maintenant introduire l'univers $\bV$, qui est d'une certaine façon
l'univers ensembliste au sens intuitif~: il est un modèle de $\ZFC$ présent dans
tout modèle de $\ZFC - \AxF$.

Historiquement, même si le nom de cet univers est attribué à Von Neumann,
il semble (d'après \cite{moore2012zermelo}) que celui-ci date des travaux de 
Zermelo, notamment \cite{Zermelo1930}.

\begin{definition}[L'univers bien fondé]
  On définit la hiérarchie d'ensembles $\bVa$, pour $\varOa \in \Ord$,
  par induction transfinie~:
  \begin{itemize}
  \item $\bV_0 \defeq \vide$.
  \item $\bV_{\varOa + 1} = \powerset(\bVa)$
  \item pour tout ordinal limite $\varOl$,
    $\displaystyle\bV_\varOl = \bigcup_{\varOb < \varOl}\bV_\varOb$
  \end{itemize}

  On définit maintenant la classe $\bV$ par
  \[\bV \defeq \bigcup_{\varOa \in \Ord} \bVa\]
\end{definition}

\begin{remark}
  Par une union indicée par $\varOa \in \Ord$, il faut en fait comprendre que le
  prédicat $\bV(x)$ est défini comme
  \[\bV(x) \defeq \exists \varOa \in \Ord, x \in \bVa\]
  On préfère une approche sémantique que syntaxique ici, car celle-ci est bien
  plus facile à appréhender.
\end{remark}

Un premier résultat important~: chaque $\bVa$ est transitif, et $\bV$ tout
entier l'est donc aussi.

\begin{property}
  Pour chaque $\varOa\in\Ord$, $\bVa$ est un ensemble transitif. De plus, $\bV$
  est une classe transitive.
\end{property}

\begin{proof}
  On raisonne par induction transfinie~:
  \begin{itemize}
  \item si $x\in \vide$, alors $x\subseteq\vide$.
  \item supposons que $\bVa$ est transitif. Soit $x\in\bV_{\varOa+1}$. Par
    définition, cela signifie que $x\subseteq \bVa$. Soit $y\in x$~: par
    inclusion, $y\in \bVa$, donc (comme $\bVa$ est transitif) $y\subseteq \bVa$,
    c'est-à-dire $y\in \bV_{\varOa + 1}$, donc $x\subseteq \bV_{\varOa + 1}$.
  \item si pour tous $\varOb < \varOl$, $\bVa$ est transitif, alors
    pour $x\in \bV_\varOl$, on trouve $\varOa < \varOl$ tel que $x\in \bVa$,
    donc $x\subseteq \bVa$. De plus, $\bVa \subseteq \bV_\varOl$ donc
    $x\subseteq \bV_\varOl$.
  \end{itemize}
  Ainsi tous les $\bVa$ sont transitifs. Le fait que $\bV$ est une classe
  transitive est une conséquence directe.
\end{proof}

Une conséquence première est que la hiérarchie est cumulative, ce qui signifie
que les ensembles sont de plus en plus grands.

\begin{property}
  Si $\varOa \leq \varOb$, alors $\bVa\subseteq \bV_\varOb$.
\end{property}

\begin{proof}
  On prouve ce résultat par induction sur l'ordinal $\varOc$ tel que
  $\varOa + \varOc = \varOb$~:
  \begin{itemize}
  \item pour tout $\varOa$, $\bVa \subseteq \bV_\varOa$.
  \item supposons que $\bVa \subseteq \bV_\varOb$, alors montrons que
    $\bVa \subseteq \bV_{\varOb+1}$~: il nous suffit pour cela de montrer que
    pour tout $\varOb$, $\bV_\varOb\subseteq\bV_{\varOb+1}$. Si $x\in\bV_\varOb$,
    alors comme $\bV_\varOb$ est transitif, $x\subseteq \bV_\varOb$, donc
    $x\in \bV_{\varOb + 1}$ par définition.
  \item supposons que pour tout $\varOd < \varOc$ et $\varOc$ ordinal limite,
    $\bVa \subseteq \bV_{\varOa + \varOd}$. Alors
    comme
    \[\bV_{\varOa + \varOc} =
    \bigcup_{\varOd < \varOc} \bV_{\varOa + \varOd}\]
    on en déduit que $\bVa \subseteq \bV_{\varOa + \varOc}$.\qedhere
  \end{itemize}
\end{proof}

Notre premier objectif est de montrer que $\bV$ constitue un premier modèle
classe de $\ZFC$. Pour cela, on donne d'abord un lemme essentiel pour prouver de
tels axiomes.

\begin{lemma}\label{lem.axZFC}
  Soit $\varCM$ une classe, alors on a les propriétés suivantes~:
  \begin{itemize}
  \item si $\varCM$ est une classe transitive, alors $\varCM$ vérifie l'axiome
    d'extensionalité.
  \item si $\varCM$ est clos par sous-ensemble, c'est-à-dire si pour tout
    ensemble $x \in \varCM$ et $y \subseteq x$, on a $y \in \varCM$, alors
    $\varCM$ vérifie l'axiome de compréhension.
  \item si pour tous $x,y \in \varCM$, $\{x,y\} \in \varCM$ alors $\varCM$
    vérifie l'axiome de la paire.
  \item si pour tout $x \in \varCM, \bigcup x \in \varCM$ alors $\varCM$ vérifie
    l'axiome de la réunion.
  \item si $\varCM$ est une classe transitive vérifiant la proposition
    suivante~:
    \begin{quotation}
      \noindent
      Pour toute fonction partielle $f$ (non obligatoirement dans $\varCM$)
      vérifiant à la fois $\dom(f) \in \varCM$ et $\im(f) \subseteq \varCM$,
      on a $\im(f) \in \varCM$.
    \end{quotation}
    alors $\varCM$ vérifie l'axiome de remplacement.
  \end{itemize}
\end{lemma}

\begin{proof}
  On prouve chaque propriété~:
  \begin{itemize}
  \item soient $x,y \in \varCM$ deux ensembles tels que
    $\forall z \in \varCM, z \in x \iff z \in y$. Alors pour tout $z \in x$,
    comme $\varCM$ est transitive, $z \in \varCM$, donc par notre hypothèse
    $z \in y$, donc $x \subseteq y$. L'argument fonctionne symétriquement pour
    nous assurer que $y \subseteq x$, donc $x = y$.
  \item supposons que $\varCM$ est clos par sous-ensemble. Soient
    $a_1,\ldots,a_n,b$ (on notera $\vec a$ pour $a_1,\ldots,a_n$)
    des éléments de $\varCM$ et $\varFP(\vec x,y)$ une formule sur
    $\SignP{\ZF}$. Alors on sait que
    $\compr{y \in b}{\varFP(\vec a,y)} \subseteq b \in \varCM$,
    donc il existe bien dans $\varCM$ un ensemble $A$ tel que
    $\forall y, y \in A \iff y \in b \land \varFP^{\varCM}(\vec x,y)$.
  \item soient $x,y \in \varCM$, alors $\{x,y\} \in \varCM$, donc il existe dans
    $\varCM$ un ensemble $A$ tel que $z \in A \iff z = x \lor z = y$.
  \item de même, si $\varCM$ est stable par réunion, on a directement un
    candidat pour être l'union d'un ensemble $x$, donné par $\bigcup x$
    lui-même~: on voit qu'un élément $z$ est dans $\bigcup x$ si et seulement
    s'il existe $y$ tel que $z \in y \in x$.
  \item soit une formule fonctionnelle $\varFP^\varCM(\vec z,x,y)$,
    $\vec a,b \in \varCM$ des paramètres. Par l'axiome de remplacement (dans
    l'univers ambiant), on sait qu'il existe un ensemble $c$, image de $a$
    par $\varFP^\varCM(\vec a)$, et une fonction $f : b \partialto c$.
    Alors, par le \cref{lem.relat.deduc}, pour tout $y \in c$, comme on sait que
    $\varFP^\varCM(\vec a,x,y)$ pour un certain $x \in b$ (donc $x \in \varCM$
    d'après le fait que $\varCM$ est transitif), on en déduit que $y\in\varCM$.
    Ainsi, $f : b \partialto c$, dont le domaine est $a \in \varCM$, a une image
    $c \subseteq \varCM$. Par notre hypothèse, $c \in \varCM$, donc l'axiome de
    remplacement est vérifié.\qedhere
  \end{itemize}
\end{proof}

\begin{remark}
  Si $\varCM$ est une classe transitive, alors elle vérifie l'axiome
  d'extensionalité mais aussi l'union, puisque $\bigcup x \subseteq \trcl(x)$.
\end{remark}

\begin{theorem}\label{thm.V.ZFCAF}
  Pour tout $\varFF \in \ZFC$, on a $\ZFC - \AxF \vdP \varFF^\bV$.
\end{theorem}

\begin{proof}
  On vérifie chaque axiome relativisé à l'aide du \cref{lem.axZFC}~:
  \begin{itemize}
  \item $\bV$ est transitive, donc $\bV$ vérifie l'extensionalité et la réunion.
  \item vérifions l'axiome de l'ensemble des parties~: soit $x\in \bV$,
    on trouve $\varOa$ tel que $x \in\bVa$. Alors pour tout $y\subseteq x$,
    $y\subseteq \bVa$ donc $y\in\bV_{\varOa +1}$. On en déduit que pour tout
    $y\in \powerset(x)$, $y\in \bV$, donc $\powerset(x)\in\bV$.
  \item vérifions l'axiome de remplacement~: soit $f$ une fonction partielle
    dont le domaine est un ensemble $x \in \bVa$, et dont l'image est
    $y \subseteq \bV$. On peut alors considérer pour chaque $z \in y$ un
    $\varOb_z$ tel que $z \in \bV_{\varOb_z}$. En prenant
    $\varOb = \sup_{z\in y} \varOb_z$, on a $y \subseteq \bV_\varOb$, donc
    $y\in \bV_{\varOb + 1}$ par définition, donc $y \in \bV$.
  \item vérifions l'axiome de l'infini~: on sait que $\Ord\subseteq\bV$, donc il
    existe bien un ensemble $\Ow^\bV \in \bV$
    contenant $0$ et stable par successeur, c'est $\Ow$ lui-même.
  \item vérifions l'axiome de fondation~: par l'absurde, supposons qu'il existe
    un élément $x \in \bV$ non vide tel que pour tout $y \in x$
    $x \cap y \neq \vide$. Soit $\varOa$ l'ordinal minimal tel que
    $x \in \bV_\varOa$. On sait que pour tout $y \in x$, il existe
    $\varOb < \varOa$ tel que $y \in \bV_\varOb$~: prenons le $\varOb$ minimal
    contenant au moins un $y \in x$ et $y$ un tel élément de $x$ associé. On
    sait alors, par hypothèse, qu'il existe $z \in x \cap y$. On en déduit,
    comme $z \in y$, qu'il existe $\varOc < \varOb$ tel que
    $z \in \bV_\varOc$, mais comme $z \in x$, cela contredit la minimalité
    de $\varOb$. C'est donc absurde, donc $x$ est bien fondé pour $\in$.
  \item vérifions l'axiome du choix~: soit un ensemble $X$ non vide, dont les
    éléments sont non vides et deux à deux disjoints. Soit $\varOa$ tel que
    $X \in \bVa$, on sait donc que tous les éléments de $X$ sont
    des éléments de $\bVa$. Par l'axiome du choix, on trouve
    $\varCh$ tel que $\varCh \cap x$ est un singleton pour tout $x \in X$.
    Quitte à restreindre $\varCh$ aux éléments d'éléments de $X$, on a donc un
    ensemble dont les éléments sont des éléments de $\bVa$, donc
    $\varCh\in\bV_{\varOa + 1}$, d'où le réstultat.\qedhere
  \end{itemize}
\end{proof}

On a donc notre premier modèle classe de $\ZFC$. Un point important lorsque l'on
ne travaille pas \latinexpr{a priori} avec l'axiome de fondation est que la
classe $\bV$ correspond exactement à la classe des éléments bien fondés de
l'univers ambiant.

\begin{proposition}
  Soit un ensemble $x$, alors on a l'équivalence suivante dans $\ZFC-\AxF$~:
  \[x\in \bV \iff
  \forall y \in \trcl(x), \exists z \in y, y \cap z = \vide\]
\end{proposition}

\begin{proof}
  L'implication directe est une conséquence du fait que $\bV\models \AxF$.
  Pour le sens réciproque, on raisonne par induction bien fondée sur $\in$~:
  si pour tout $x \in X$, $x \in \bV$, alors en considérant
  \[\rk(x) \defeq \min\compr{\varOa \in \Ord}{x \subseteq\bV_\varOa}\]
  puis
  \[\varOa \defeq \sup_{x \in X} \rk(x)\]
  on sait que $x\in \bV_{\varOa+1}$ pour tout $x \in X$, donc
  $X\subseteq \bV_{\varOa+1}$, donc $X \in \bV_{\varOa + 2}$.
  Ainsi, pour tout ensemble $X$ tel que $\in$ est bien fondé sur $X$,
  $X \in \bV$.
\end{proof}

On voit donc que dans le cas d'un univers ensembliste $\unZF$, celui-ci
vérifie $\AxF$ si et seulement si $\bV = \unZF$. On utilisera donc en
général $\bV$ pour parler de l'univers ensembliste.

\section{Vers le théorème de réflexion}

Dans le reste du chapitre, on suppose l'axiome de fondation~: on considère donc
que $\unZF = \bV$. Notre premier objectif est de montrer le théorème
de réflexion. Pour ce faire, nous introduisons d'abord la notion de rang,
permettant d'étudier plus en détail la hiérarchie des $\bV_\varOa$, et les
ensembles $\bH_\varCk$. Nous verrons ensuite l'absoluité, permettant de rendre
plus formel ce qu'on attend d'une proposition qui passe d'un modèle de $\ZFC$
à un autre.

\subsection{Rang et parties héréditaires}

Le rang indique à quelle étape de la construction des $\bV_\varOa$ notre
ensemble a été construit.

\begin{definition}[Rang]
  On définit le rang d'un ensemble $x\in\bV$ par~:
  \[\rk(x) \defeq \min \compr{\varOa \in \Ord}{x\subseteq \bVa}\]
\end{definition}

Le rang d'un ensemble peut se voir comme la profondeur d'accolades imbriquées
lorsque l'on définit l'ensemble en question. Ce rang se comporte bien vis à vis
des opérations ensemblistes usuelles, comme nous allons le voir.

Pour mieux étudier la relation entre les opérations ensemblistes et le rang,
on introduit une caractérisation du rang, qui pourrait être vue comme une
définie par récursion sur $\in$ (récursion bien fondée dans le cas présent
puisqu'on suppose $\AxF$).

\begin{proposition}\label{prop.cara.rk}
  Soit un ensemble $X\in\bV$, on peut caractériser son rang par
  \[\rk(X) = \sup_{x \in X} (\rk(x) + 1)\]
\end{proposition}

\begin{proof}
  Montrons que $\rk(X)$ est un majorant de $\compr{\rk(x)+1}{x \in X}$.
  Pour cela, on montre que pour tout $x \in X$, si $X\subseteq \bVa$
  alors $x\subseteq \bVa$ (donc le minimum des tels $\bVa$ contenant $x$ est
  inférieur au minimum de ceux contenant $X$). Soit donc $x \in X$ et $\varOa$
  tel que $X\subseteq \bVa$, alors $x\in \bVa$, mais comme $\bVa$ est transitif,
  on en déduit que $x\subseteq \bVa$. Donc $\rk(X)$ est un majorant de
  $\compr{\rk(x) + 1}{x \in X}$.

  Montrons que ce majorant est une borne supérieure. Soit $\varOb\in\Ord$ un
  majorant de $\compr{\rk(x)+1}{x \in X}$. Pour tout $x\in X$,
  $x\subseteq \bV_{\rk(x)}$ donc $x\in \bV_{\rk(x)+1}$ donc $X\subseteq \bV_\varOb$,
  donc $\rk(X) \leq \varOb$. Ce majorant est donc une borne supérieure, d'où
  l'égalité.
\end{proof}

En voyant le rang comme le nombre d'accolades imbriquées, on peut remarquer
que chaque rang inférieur à un $\rk(x)$ donné est présent dans $x$. Non
forcément en tant que rang d'un de ses éléments, mais potentiellement d'un
élément d'un de ses éléments \latinexpr{etc.} car, s'il y avait un \og saut\fg
de rang dans les éléments de $x$, alors les éléments au-dessus de ce saut
n'auraient pas de raison d'être déjà présents plus tôt dans la hiérarchie des
$\bVa$.

\begin{lemma}\label{lem.rk.surj}
  Soit $x$ un ensemble et $\varOa$ son rang. Alors pour tout $\varOb < \varOa$,
  il existe $y$ tel que $\varOb = \rk(y)$ et $y \in \trcl(x)$.
\end{lemma}

\begin{proof}
  Par l'absurde, en fixant $\varOa$, supposons qu'il existe un ensemble $x$,
  minimal pour $\in$, tel que $\rk(x) > \varOa$ et il n'existe aucun élément
  $y\in\trcl(x)$ tel que $\rk(y) = \varOa$. Comme on sait que $\varOa < \rk(x)$
  et que $\displaystyle \rk(x) = \sup_{y \in x} (\rk(y) + 1)$,
  alors on peut trouver $y \in x$ tel que $\rk(y) + 1 > \varOa$. Si
  $\rk(y) = \varOa$, alors cela contredit notre hypothèse sur $x$. Si
  $\rk(y) > \varOa$, alors par $\in$-minimalité de $x$, on en déduit qu'il
  existe $z \in \trcl(y)$ tel que $\rk(z) = \varOa$, mais alors
  $z \in \trcl(x)$, ce qui est absurde. Ainsi, dans tous les cas, on a une
  absurdité. On en déduit que pour tout $x$ et $\varOa < \rk(x)$, il existe
  $y \in \trcl(x)$ tel que $\rk(y) = \varOa$.
\end{proof}

\begin{property}\label{prop.rg.oper}
  Les propriétés suivantes sont vérifiées~:
  \begin{itemize}
  \item $\forall x,y\in\bV, x\in y \implies \rk(x) + 1 \leq \rk(y)$
  \item $\forall x,y\in\bV, x\subseteq y \implies \rk(x) \leq \rk(y)$
  \item $\displaystyle\forall x\in\bV,\rk\left(\bigcup x\right)+1\leq\rk(x)$
  \item $\forall x,y\in\bV, \rk(x\cup y) = \max(\rk(x),\rk(y))$
  \item $\forall x,y\in\bV, \rk(\{x,y\}) = \max(\rk(x),\rk(y)) + 1$
  \item $\forall x\in\bV, \rk(\powerset(x)) = \rk(x) + 1$
  \item $\forall x\in\bV, \rk(\trcl(x)) = \rk(x)$
  \item $\displaystyle \forall \varOa \in \Ord, \rk(\varOa) = \varOa$
  \end{itemize}
\end{property}

\begin{proof}
  Vérifions chaque propriété~:
  \begin{itemize}
  \item l'inégalité vient directement de la \cref{prop.cara.rk}.
  \item soient $x,y$ tels que $x\subseteq y$. Si $y\subseteq\bVa$, comme
    $x\subseteq y$, on en déduit que $x\subseteq\bVa$, donc $\rk(x)\leq \rk(y)$.
  \item soit $y\in \bigcup x$, on trouve $z\in x$ tel que $y\in z \in x$. En
    utilisant l'inégalité précédente, on en déduit que $\rk(y) + 1 \leq \rk(z)$
    et que $\rk(z) + 1 \leq \rk(x)$, donc $\rk(y) + 2 \leq \rk(x)$. On en déduit
    en utilisant la \cref{prop.cara.rk} que $\rk(\bigcup x) + 1 \leq \rk(x)$.
  \item soient $x,y$. On utilise directement la \cref{prop.cara.rk}~:
    \begin{align*}
      \rk(x\cup y) &= \sup_{z \in x \cup y} (\rk(z) + 1)\\
      &= \max(\sup_{z \in x} (\rk(z) + 1), \sup_{z\in y} (\rk(z) + 1))\\
      &= \max(\rk(x),\rk(y))
    \end{align*}
  \item pour deux éléments $x,y$, il est clair que le sup est un max, d'où
    l'égalité.
  \item soit un ensemble $x$. comme $x\in \powerset(x)$, on en déduit que
    $\rk(x) + 1 \leq \rk (\powerset (x))$. Pour l'autre inclusion, si
    $y\in \powerset(x)$ alors $y\subseteq x$, donc
    $\rk(y) \leq \rk(x)$ par un point précédent. On en déduit que
    $\sup_{y\in \powerset(x)} (\rk(y) + 1) \leq \rk(x) + 1$ d'où
    l'égalité.
  \item comme chaque $\bVa$ est transitif, si $x\subseteq \bVa$ alors
    $\trcl(x)\subseteq \bVa$. On en déduit que $\rk(\trcl(x)) \leq \rk(x)$, et
    il est clair que l'autre inégalité tient toujours puisqu'on rajoute des
    éléments.
  \item on prouve cela par induction transfinie~:
    \begin{itemize}
    \item $\vide \subseteq \vide$ donc $\rk(0) = 0$.
    \item on sait que $\varOa + 1 = \varOa \cup \{\varOa\}$. En utilisant les
      précédentes égalités, $\rk(\{\varOa\}) = \rk(\varOa) + 1$ donc
      $\rk(\varOa + 1) = \varOa + 1$.
    \item si $\varOl$ est un ordinal limite, alors
      $\rk(\varOl) = \sup_{\varOb < \varOl} (\rk(\varOb) + 1)$ et par hypotèse
      d'induction, $\rk(\varOb) = \varOb$, donc $\rk(\varOl) = \varOl$.
    \end{itemize}
    Ainsi pour tout $\varOa \in \Ord, \rk(\varOa) = \varOa$.\qed
  \end{itemize}
\end{proof}

On peut alors affiner certains résultats liés au fait que
$\bV \satisf \ZFC$. En effet, puisque $\bVa \subseteq\bV$, on peut chercher ce
que chaque $\bVa$ vérifie. Par exemple, chaque $\bVa$ vérifie l'axiome de
fondation.

On voit dans la \cref{prop.rg.oper} que la plupart des opérations ne font pas
exploser le rang. Cela signifie que, dans le cas d'un ordinal $\varOl$ limite,
où $\varOa \in \varOl \implies \varOa + 1 \in \varOl$, beaucoup de constructions
sont stables par $\bV_{\varOl}$.

\begin{property}
  Soit $\varOl$ un ordinal limite. Alors $\bV_\varOl$ vérifie les
  axiomes suivants~:
  \begin{itemize}
  \item l'axiome de la paire
  \item l'axiome de l'union
  \item l'axiome de l'ensemble des parties
  \item l'axiome de compréhension
  \item l'axiome du choix
  \end{itemize}

  De plus, si $\varOl > \Ow$ alors $\bV_\varOl$ vérifie l'axiome de l'infini.
\end{property}

\begin{proof}
  Vérifions chaque axiome, en utilisant le \cref{lem.axZFC}~:
  \begin{itemize}
  \item axiome de la paire~: soient $x,y\in \bV_\varOl$, donc tels que
    $\rk(x) <\varOl$ et $\rk(y)<\varOl$. On a alors l'inégalité
    $\rk(\{x,y\}) = \max(\rk(x),\rk(y)) + 1 < \varOl$, donc
    $\{x,y\}\in\bV_\varOl$.
  \item axiome de l'union~: soit $x\in \bV_\varOl$, alors $\rk(x) < \lambda$,
    donc $\rk(\bigcup x) < \varOl$ donc $\bigcup x \in \bV_\varOl$.
  \item axiome de l'ensemble des parties~: soit $x\in \bV_\varOl$. Comme
    $\rk(\powerset(x)) = \rk(x) + 1$, que $\rk(x) < \varOl$ et que $\varOl$
    est limite, on en déduit que $\rk(\powerset(x)) < \varOl$ donc que
    $\powerset(x) \in \bV_\varOl$.
  \item axiome de compréhension~: on sait que si $x\subseteq y$ alors
    $\rk(x) \leq \rk(y)$, donc $\bV_{\varOl}$ est stable par
    sous-ensemble.
  \item soit $X$ un ensemble non vide, d'éléments deux à deux disjoints. Alors
    par l'axiome du choix, il existe $\varCh$ tel que
    $\forall x \in X,\Card(\varCh\cap x)=1$.
    Quitte à prendre $\varCh \cap \bigcup X$, $\varCh\subseteq \bigcup X$ donc
    $\rk(\varCh) < \varOl$, donc $\varCh\in \bV_\varOl$.
  \end{itemize}

  Si $\varOl > \Ow$, alors $\Ow \in \bV_\varOl$ donc il existe bien dans
  $\bV_\varOl$ un élément contenant l'ensemble vide et stable par successeur. On
  peut préciser que l'ensemble vide est bien dans $\bV_\varOl$, donnant un sens à
  la définition.
\end{proof}

On peut se demander alors si l'axiome de remplacement est vérifié. En fait il
est heureux que ça ne soit pas le cas, en vertu du deuxième théorème
d'incomplétude de Gödel~: si c'était le cas, alors on trouverait un modèle de
$\ZFC$, impliquant que $\ZFC\vdP \Coher(\ZFC)$ (ou plutôt
$\ZFC\vdP \Coher(\godcod{\ZFC})$).

Les $\bV_\varOl$ nous donnent donc une approximation assez fidèle d'un modèle de
$\ZFC$ s'il s'agit de mathématiques usuelles, mais ces ensembles perdent de leur
intérêt pour de la théorie des ensembles~: le schéma d'axiomes de remplacement
est utilisé pour la plupart de nos constructions, dès l'utilisation de la
récursion transfinie.

On introduit une autre hiérarchie d'ensembles qui, eux, vérifieront le schéma
d'axiomes de remplacement~: ce sont les ensembles héréditaires.

\begin{definition}[Ensemble héréditaire]
  Soit $\varCk \in \Card$. On appelle ensemble des ensembles héréditairement de
  cardinal $\varCk$ la classe
  \[\bH_\varCk \defeq \compr{x}{\Card(\trcl(x)) < \varCk}\]
\end{definition}

Pour deux cardinaux $\varCk,\varCCl$, on définit le cardinal
\[\varCCl^{<\varCk} \defeq \sup_{\varCk' < \varCk} \varCCl^{\varCk'}\]
On peut voir ce cardinal comme le cardinal de l'ensemble des fonctions de
$\varCk$ dans $\varCCl$ dont le support (ensemble des points de $\varCk$ où la
fonction n'est pas nulle) est strictement inférieur à $\varCk$.

Un premier point essentiel sur cette hiérarchie~: les $\bH_\varCk$ sont des
ensembles. Pour le montrer, on se repose sur une inclusion de $\bH$ dans la
hiérarchie des $\bV$, due au fait que tout rang est représenté dans la clôture
transitive d'un ensemble.

\begin{proposition}\label{prop.Hcard}
  Soit $\varCk\in\Card$ infini. Alors $\bH_\varCk\subseteq\bV_\varCk$
  et $\Card(\bH_\varCk) = 2^{<\varCk}$.
\end{proposition}

\begin{proof}
  Soit un ensemble $x$ et $\varOa$ son rang. Par le \cref{lem.rk.surj}, pour
  tout $\varOb < \varOa$, on trouve $x_\varOb \in \trcl(x)$ tel que
  $\rk(x_\varOb) = \varOb$. On a donc une injection de $\rk(x)$ vers $\trcl(x)$,
  donc $\Card(\trcl(x)) \geq \Card(\rk(x))$, d'où $x \in \bV_\varCk$.

  Pour la deuxième proposition, procédons par double inégalité. Tout d'abord,
  pour $\varCCl < \varCk$, on sait que $\powerset(\varCCl) \subseteq \bH_\varCk$,
  donc $\Card(\bH_\varCk) \geq 2^\varCCl$, d'où $\Card(\bH_\varCk) \geq 2^{<\varCk}$.

  Pour l'autre inégalité, on va construire pour chaque $x \in \bH_\varCk$
  une relation $\varRR_x \subseteq \varCCl \times \varCCl$ pour un certain
  $\varCCl < \varCk$, de telle sorte que l'association $x \mapsto \varRR_x$ soit
  injective. Soit donc $x \in \bH_\varCk$, on pose
  $\varCCl = \Card(\trcl(x)\cup \{x\})$ (donc $\varCCl < \varCk$) et on pose
  $\varRR_x\subseteq \varCCl \times \varCCl$ une relation telle que
  \[(\varCCl,\varRR_x)\cong (\trcl(x)\cup \{x\}, \in)\]
  (cette relation existe puisqu'on a une bijection entre $\trcl(x)\cup\{x\}$ et
  $\varCCl$). On veut maintenant prouver que pour tout autre $y$ et
  $\varCCl' = \Card(\trcl(y)\cup\{y\})$, si
  $(\varCCl, \varRR_x) = (\varCCl',\varRR_y)$, alors $x = y$. On suppose donc
  que $(\varCCl, \varRR_x) = (\varCCl',\varRR_y)$. Par le
  \cref{lem.most}, on sait (puisque $\trcl(x)\cup \{x\}$ et $\trcl(x)\cup \{y\}$
  sont des ensembles transitifs) que le seul isomorphisme existant entre
  $(\trcl(x)\cup \{x\},\in)$ et $(\trcl(y)\cup\{y\},\in)$ est l'identité. Or,
  on sait que $(\lambda,r_x) = (\lambda',r_y)$, donc on en conclut que
  $\trcl(x)\cup \{x\} = \trcl(y)\cup \{y\}$ par unicité de l'isomorphisme.
  Il vient ensuite que $x = y$ puisque $x$ (respectivement $y$) est l'unique
  élément $\in$-maximal dans $\trcl(x)\cup \{x\}$ (respectivement dans
  $\trcl(y)\cup\{y\}$). On a donc une injection de $\bH_\varCk$ dans
  $\displaystyle\bigcup_{\varCCl < \varCk}\lambda \times 2^{\varCCl \times \varCCl}$,
  d'où $\Card(\bH_\varCk) \leq 2^{<\varCk}$.
\end{proof}

\noindent
Comme $\bH_\varCk$ est inclus dans un ensemble, c'est lui-même un ensemble.

On veut maintenant montrer les cas d'égalité entre $\mathbb H$ et $\mathbb V$.

\begin{property}
  Les propriétés suivantes sont vérifiées~:
  \begin{itemize}
  \item $\bV_\Ow = \bH_\Ow$
  \item pour tout $\varOa\in\Ord$, $\Card(\bV_{\Ow + \varOa}) = \Cbe_\varOa$,
    donc en particulier pour $\varOa \geq \Ow^2$, $\Card(\bVa) = \Cbe_\varOa$.
  \item si $\varCk > \Ow$, alors $\bH_\varCk = \bV_\varCk$ si et seulement si
    $\varCk = \Cbe_\varCk$.
  \end{itemize}
\end{property}

\begin{proof}
  Prouvons chaque propriété~:
  \begin{itemize}
  \item on voit que pour tout $n$, $\bV_n$ ne contient que des ensembles finis
    (c'est le cas de $\bV_0$, et $\powerset(X)$ ne contient que des ensembles
    finis si $X$ est fini), donc $\bV_n \subseteq \bH_\Ow$, d'où
    $\bV_\Ow \subseteq \BH_\Ow$. Réciproquement, si $x \notin \bV_\Ow$, alors
    $\trcl(x)$ est au moins dénombrable, puisque $x$ a un élément de chaque
    rang $n \in \Ow$~: par contraposée, si $x\in \bH_\Ow$ alors $x\in \bV_\Ow$.
  \item on sait que $\Card(\bV_\Ow) = \Ow$ car on a une union dénombrable
    d'ensembles finis, que $\Card(\bV_{\varOa + 1}) = 2^{\bV_\varOa}$ et que
    $\Card(\bV_\varOl) = \sup_{\varOb < \varOl} \Card(\bV_\varOb)$ pour
    $\varOl$ limite, donc $\bV_{\Ow + \varOa}$ coïncide avec $\Cbe_\varOa$. Si
    $\varOa \geq \Ow^2$, alors $\Ow + \varOa = \varOa$, d'où le résultat.
  \item on prouve l'équivalence par double implication~:
    \begin{itemize}
    \item soit $\varCk > \Ow$ tel que $\bH_\varCk = \bV_\varCk$.
      Soit $\varOa$ tel que $\Ow^2\leq \varOa < \varCk$. On sait que
      $\bVa \in \bV_\varCk$, donc $\bVa \in \bH_\varCk$, donc en utilisant la
      propriété prouvée précédemment, $\Cbe_\varOa < \varCk$. En prenant alors la
      borne supérieure de ces égalités pour $\varOa < \varCk$ et par continuité
      de $\Cbe$, on en déduit que $\Cbe_\varCk \leq \varCk$. Comme on sait que
      $\varCk \leq \Cbe_\varCk$ pour n'importe quel $\varCk$, on en déduit que
      $\varCk = \Cbe_\varCk$.
    \item supposons que $\varCk = \Cbe_\varCk$. Soit $x \in \bV_\varCk$,
      on peut trouver $\varOa$ tel que $\Ow^2\leq \varOa < \varCk$ et
      $x \in \bV_\varOa$, car $\varCk$ est un ordinal limite et un cardinal
      indénombrable. Comme $x \subseteq \bV_\varOa$, on en déduit que
      $\Card(\trcl(x)) \leq \Card(\bVa)$, or $\bVa$ est
      de cardinal $\Cbe_\varOa$, et $\Cbe_\varOa < \Cbe_\varCk = \varCk$. On en
      déduit donc que $x \in \bH_\varCk$. L'inclusion réciproque est directement
      due à la \cref{prop.Hcard}.\qedhere
    \end{itemize}
  \end{itemize}
\end{proof}

On peut maintenant prouver que $\bH_\varCk$ vérifie les axiomes de
$\ZFC - \AxPow$ pour un cardinal régulier $\varCk$, c'est-à-dire les axiomes de
$\ZFC$ sauf celui de l'ensemble des parties.

\begin{property}
  Soit $\varCk$ un cardinal régulier indénombrable. Alors $\bH_\varCk$
  vérifie~:
  \begin{itemize}
    \begin{minipage}{0.4\textwidth}
    \item l'axiome d'extensionalité
    \item l'axiome de la paire
    \item l'axiome de la réunion
    \end{minipage}
    \begin{minipage}{0.4\textwidth}
    \item l'axiome de remplacement
    \item l'axiome de l'infini
    \item l'axiome de fondation
    \end{minipage}
  \item l'axiome du choix
  \end{itemize}
\end{property}

\begin{proof}
  On vérifie chaque axiome, en utilisant le \cref{lem.axZFC}~:
  \begin{itemize}
  \item axiome de la paire~: comme
    $\trcl(\{x,y\}) = \trcl(x) \cup \trcl(y) \cup \{x,y\}$, si
    $\trcl(x)$ et $\trcl(y)$ sont de cardinal inférieur strictement à $\varCk$,
    alors $\Card(\trcl(\{x,y\})) < \kappa$.
  \item axiome de la réunion~: si $X\in \bH_\varCk$, alors
    $\Card(\trcl(X)) < \varCk$. Or, on peut remarquer que
    \[\trcl\left(\bigcup X\right)\subseteq \trcl(X)\]
    puisque $\trcl(X)$ contient déjà $\bigcup X$ et est transitif. Il vient donc
    que $\displaystyle\Card\left(\trcl\left(\bigcup X\right)\right) < \varCk$.
  \item axiome de remplacement~: soit $f$ une fonction partielle d'un
    ensemble $x \in \bH_\varCk$ vers un ensemble $y \subseteq \bH_\varCk$. Alors
    \[\trcl(y) = \bigcup_{z \in y} (\trcl(z)\cup \{z\})\]
    mais comme chaque $z \in y$ est un élément de $\bH_\varCk$, et qu'on sait que
    $x \in \bH_\varCk$, on en déduit que $\trcl(y)$ est une union de cardinal
    inférieur strictement à $\varCk$ d'ensembles de cardinal inférieur
    strictement à $\varCk$. Ainsi, d'après le \cref{thm.Konig},
    $|\trcl(y)| < \varCk$, donc $y \in \bH_\varCk$.
  \item axiome de l'infini~: puisque $\varCk$ est indénombrable, il contient
    $\Ow$.
  \item axiome de fondation~: $\bH_\varCk \subseteq \bV_\varCk$ donc
    tous les éléments de $\bH_\varCk$ sont bien fondés pour $\in$.
  \item axiome du choix~: si un ensemble $x \in \bH_\varCk$ est non vide
    et tous ses éléments sont non vide et deux à deux disjoints, alors par
    l'axiome du choix on trouve $\varCh$ tel que $\varCh \cap y$ est un
    singleton pour tout $y \in x$. Quitte à intersecter $\varCh$ avec
    $\bigcup x$, on sait que $\varCh\subseteq \bigcup x$, donc comme
    $\bH_\varCk$ est stable par union et par sous-ensembles,
    $\varCh \in \bH_\varCk$.
  \end{itemize}
  Ainsi $\bH_\varCk$ vérifie $\ZF - \AxPow$.
\end{proof}

On voit qu'ici, pour que $\bH_\varCk$ vérifie tout $\ZFC$, il lui suffit que pour
tout $\varCk' \in \varCk$, $2^{\varCk'} \in \varCk$. C'est ce qu'on appelle un
cardinal fortement inaccessible.

\begin{definition}[Cardinal fortement inaccessible]
  On dit qu'un cardinal $\varCCl$ est fortement inaccessible (on dit aussi, plus
  simplement, inaccessible) si $\varCCl > \Ow$ et
  \[\forall \varCk \in \varCCl, 2^\varCk \in \varCCl\]
\end{definition}

Il vient directement qu'un tel cardinal ne peut exister dans tout modèle de
$\ZFC$.

\begin{proposition}
  Il n'est pas prouvable dans $\ZFC$ qu'il existe un cardinal inaccessible.
\end{proposition}

\begin{proof}
  S'il existait un $\varCk$ inaccessible, alors $\bH_\varCk$ serait un modèle
  ensembliste de $\ZFC$, ce qui contredit le second théorème d'incomplétude de
  Gödel.
\end{proof}

Enfin, nous allons introduire l'univers constructible $\bL$ de Gödel, que
nous étudierons dans la section prochaine. Nous aurons cependant besoin de
notions liées à sa construction.

\begin{definition}[Partie définissable]
  Soit un ensemble $X$, on dit qu'une partie $Y\subseteq X$ est définissable
  dans $X$ s'il existe une formule $\varFF(\vec x,y)$ à $n+1$ variables
  libres et des ensembles $\vec a \in X^n$ tels que
  \[\forall b \in X, b \in Y \iff X\satisf \varFF(\vec a,b)\]

  \noindent
  On appelle ensemble des parties définissables d'un ensemble $X$ l'ensemble
  %\begin{multline*}
  \[\powerdef(X) \defeq \compr{Y \subseteq X}{\exists
    \varFF(\vec x,y)\in\Formula_{n+1}(\SignP{\ZF}),
    \exists \vec a \in X^n, \forall b \in X, b \in Y \iff X \satisf
    \varFF(\vec a,b)}\]
  %\end{multline*}
\end{definition}

\begin{remark}
  Il peut être bon de préciser, ici, de quel type de formule on parle. Dans le
  cas de la notion de définissabilité, on considère les propositions internes
  au modèle de $\ZFC$, c'est-à-dire les objets ensemblistes représentant
  intérieurement nos propositions. Cela explique l'utilisation de
  \og $X\satisf \varFF(x)$\fg plutôt que \og$\varFF(x)$\fg directement.
\end{remark}

L'univers constructible est alors l'itération de cette opérations à la place de
l'ensemble des parties.

\begin{definition}[Univers constructible de Gödel\cite{GödelL}]\label{def.L}
  On définit par induction transfinie la hiérarchie $\mathbb L_\alpha$~:
  \begin{itemize}
  \item $\bL_0 = \vide$
  \item $\bL_{\varOa + 1} = \powerdef (\bL_\varOa)$
  \item si $\varCCl$ est limite, alors
    $\displaystyle\bL_\varOl = \bigcup_{\varOb < \varOl} \bL_\varOb$
  \end{itemize}
  L'univers constructible de Gödel est alors
  \[\bL \defeq \bigcup_{\varOa \in \Ord} \bL_\varOa\]
\end{definition}

\subsection{Absoluité et réflexion}

On va maintenant s'intéresser à une question assez naturelle à se poser, dans
les conditions de notre étude des modèles de $\ZFC$. Supposons donnés deux
modèles (potentiellement des classes) $(\varCM,\in)$ et $(\varCM',\in)$
et $\varCM \subseteq \varCM'$, quand une formule $\varFF$ à paramètres dans
$\varCM$ garde-t-elle la même valeur de vérité que l'on l'étudie dans $\varCM$
ou dans $\varCM'$~?

Prenons un exemple~: on a vu que $\bVa$ étant un ensemble transitif, celui-ci
vérifiait l'extensionalité. La preuve montre assez clairement que la proposition
$\forall z, z \in x \iff z \in y$ en elle-même ne dépend pas du point de vue
choisi, tant que celui-ci se fait dans un ensemble transitif. Ainsi si on
imagine un sur-modèle de notre univers ensembliste, ou une partie seulement, la
vérité de $\forall z, z \in x \iff z \in y$ est presque constante (encore une
fois, tant que l'univers considéré est transitif). Une telle formule est
absolue~: sa valeur de vérité ne change pas suivant les modèles.

\begin{definition}[Formule absolue]
  Soient deux classes $(\varCM,\in)$, $(\varCM',\in)$ telles que
  $\varCM \subseteq \varCM'$.
  Soit $\varFF(\vec x) \in \Formula_n(\SignP{\ZF})$ une formule. On dit que
  $\varFF$ est absolue entre $\varCM$ et $\varCM'$ lorsque
  \[\varCM \absolF \varCM' \defeq
  \forall \vec a \in \varCM^n, \varFF^{\varCM}(\vec a)
  \iff \varFF^{\varCM'}(\vec a)\]

  On dira que $\varFF$ est absolue pour une classe $\varCM$ si elle est
  absolue entre $\varCM$ et $\bV$.
\end{definition}

L'exemple précédent sur l'extensionalité peut en fait se généraliser. La
propriété exploitée est le fait que si l'on étudie
$\forall x, z \in x \iff z \in y$ dans un modèle transitif, alors on est obligé
de contenir tous les éléments de $x$ et $y$. Cette propriété est en quelque
sorte locale~: tout ce qui est nécessaire pour évaluer la vérité de cette
proposition est la donnée de tous les éléments de $x$ et de $y$.

Au contraire, une propriété comme \og être un cardinal\fg n'est pas du tout
locale, puisqu'elle repose sur la non existence d'une injection dans l'univers
ambiant. Imaginons qu'on puisse ajouter une telle injection dans notre univers,
alors un élément qui était un cardinal peut ne plus l'être.

Pour étudier les propositions qui présentent cette propriété de localité
mentionnée ci-dessus, on introduit la hiérarchie de Lévy.

\begin{definition}[Hiérarchie de Lévy \cite{LevyH}]\label{def.Levy}
  On définit par induction trois familles d'ensembles de formules sur
  $\SignP{\ZF}, (\HierLevD{n})_{n\in\bN}, (\HierLevS{n})_{n\in\bN}$
  et $(\HierLevP{n})_{n \in \bN}$~:
  \begin{itemize}
  \item $\DeltaZL=\HierLevS{0}=\HierLevP{0}$ est l'ensemble des formules
    prouvablement équivalentes dans $\ZFC$ à une formule dont les
    quantificateurs sont tous bornés, c'est-à-dire que toute quantification
    $\quantif x, \varFP$ est de la forme $\quantif x \in t, \varFP$ où $t$ est
    un terme ne contenant pas $x$.
  \item $\\HierLevS{n+1}$ est l'ensemble des formules prouvablement équivalentes
    dans $\ZFC$ à une formule de la forme $\exists x, \varFP$ où $\varFP$ est
    une formule $\HierLevP{n}$.
  \item $\HierLevP{n+1}$ est l'ensemble des formules prouvablement équivalentes
    dans $\ZFC$ à une formule de la forme $\forall x, \varFP$ où $\varFP$ est
    une formule $\HierLevS{n}$.
  \item $\HierLevD{n}$ est l'ensemble des formules qui sont à la fois
    $\HierLevS{n}$ et $\HierLevP{n}$.
  \end{itemize}

  %  On notera dans ce chapitre simplement $\Delta_n$ (respectivement $\Sigma_n$
  %  et $\Pi_n$) pour $\Delta^{\ZFC}_n$.
\end{definition}

\begin{remark}
  Comme on peut toujours prendre un produit fini comme un nouvel ensemble, il
  est équivalent de prendre un seul quantificateur $\exists$ (respectivement
  $\forall$) et une suite finie de quantificateurs $\exists$ (respectivement
  $\forall$), d'où le fait qu'on se limite dans la hiérarchie aux alternances
  de quantificateurs.
\end{remark}

On peut maintenant énoncer la généralisation de notre preuve de l'axiome
d'extensionalité pour toute formule $\DeltaZL$.

\begin{proposition}
  Soit $\varCM$ une classe transitive. Alors toute formule $\varFF\in\DeltaZL$
  est absolue pour $\varCM$.
\end{proposition}

\begin{proof}
  On prouve le résultat par induction sur la formule $\varFP$ qu'on considère~:
  \begin{itemize}
  \item dans le cas d'une formule atomique, puisque $\in$ est la même relation
    interprétée dans $\varCM$ et dans $\bV$, on a le résultat (de même pour
    $=$).
  \item dans le cas de $\lor,\land,\to,\lnot,\top,\bot$, le résultat se fait
    simplement en écrivant les définitions.
  \item dans le cas de $\exists$, supposons donc que
    $\varFP(\vec x,z) = \exists y \in z,\varFF(\vec x,y,z)$
    (les termes dans $\SignP{\ZF}$ sont uniquement des variables). Par hypothèse
    d'induction, on suppose donc que $\varFF$ est absolue pour $\varCM$. Soient
    des paramètres $\vec a$ et $c$ dans $\varCM$. On veut donc prouver que
    \[\exists b^{\varCM}, b \in c \land \varFF^{\varCM}(\vec a,b,c)
    \iff \exists b \in c, \varFF(\vec a,b,c)\]
    Le sens direct se vérifie facilement. Pour le sens réciproque, on remarque
    que comme $\varCM$ est transitive, $b \in c$ implique que $b \in \varCM$.
    Alors, par absoluité de $\varFF$ pour $\varCM$, on en déduit que
    $\varFF(\vec a,b,c)$.
  \item on peut prouver le cas de $\forall$ en passant par les loi de De Morgan.
  \end{itemize}
  Ainsi, par induction, on sait que $\varFP \in \DeltaZL$ est absolue pour
  $\varCM$.
\end{proof}

De nombreuses formules sont $\DeltaZL$. En particulier celles définissant
plusieurs objets essentiels de la théorie des ensembles.

\begin{example}
  Les formules suivantes sont $\DeltaZL$~:
  \begin{itemize}
  \item $x\in y$
  \item $x\subseteq y$, car cela s'écrit $\forall z \in x, z \in y$
  \item $x = \{y,z\}$
  \item $x = \bigcup y$
  \item $x$ est transitif, car cela s'écrit $\forall y \in x, y\subseteq x$
  \item $x$ est un couple, car cela s'écrit
    $\exists y,\exists z, x = \{\{y,z\},\{y\}\}$
  \item pour un couple $\alpha$, la formule $x = \proj{1}{}(\alpha)$
    (respectivement $x = \proj{2}{}(\alpha)$)
  \item $f$ est une fonction, car cela s'écrit
    \[\left(\forall \gamma \in f, \gamma\text{ est un couple }\right)\land
    \left(\forall \gamma,\delta \in f, \proj{2}{}(\gamma) = \proj{2}{}(\delta)
    \implies \proj{1}{}(\gamma) = \pi{1}{}(\delta)\right) \]
  \item $\varOa \in \Ord$, car être un ordinal est la conjonction des formules
    disant~:
    \begin{itemize}
    \item que $\varOa$ est transitif, dont on a vu que la formule était
      $\DeltaZL$.
    \item que $\in$ est transitif sur $\varOa$, ce qui s'écrit
      \[\forall x \in \varOa, \forall y \in \varOa, \forall z \in \varOa,
      x\in y \land y \in z \implies x \in z\] qui est clairement $\DeltaZL$.
    \item que $\in$ est total, ce qui s'écrit
      \[\forall x \in \varOa, \forall y \in \varOa, x \neq y \implies
      x \in y \lor y \in x\] qui est là encore clairement $\DeltaZL$.
    \item comme on suppose l'axiome de fondation, l'irréflexivité de $\in$ et sa
      bonne fondation sont forcément vérifiées.
    \end{itemize}
  \item être égal à $\Ow$, car \og $X = \Ow$\fg peut s'écrire
    \[X\in \Ord \land (\forall z \in X, \Ss z \in X) \land
    (\forall y \in X, \lnot (\forall z \in y, \Ss z \in y))\]
  \end{itemize}
\end{example}

\begin{exercise}
  Montrer qu'être un ordinal successeur (respectivement limite) peut s'exprimer
  par une formule $\DeltaZL$. Montrer qu'être une bijection peut s'exprimer
  par une formule $\DeltaZL$.
\end{exercise}

On voit donc une famille importante de formules qui sont absolues entre les
classes transitives (qui sont presque les seules classes d'intérêt).

On peut aussi étudier les formules $\HierLevS{1}$ et $\HierLevP{1}$, qui ne sont
pas absolues mais transmettent leur vérité d'une classe à une sur-classe
(respectivement à une sous-classe).

\begin{proposition}
  Soient $\varCM$, $\varCM'$ deux classes transitives, avec
  $\varCM \subseteq \varCM'$. Alors pour un énoncé
  $\varFF\in\Formula(\SignP{\ZF})$, on a les implications suivantes~:
  \begin{itemize}
  \item si $\varFF \in \HierLevS{1}$ alors
    $\varFF^{\varCM} \implies \varphi^{\varCM'}$
  \item si $\varFF \in \HierLevP{1}$ alors
    $\varFF^{\varCM'} \implies \varphi^{\varCM}$
  \end{itemize}
  et toute formule $\HierLevD{1}$ est absolue entre $\varCM$ et $\varCM'$.
\end{proposition}

\begin{proof}
  On montre d'abord les deux implications~:
  \begin{itemize}
  \item supposons que $\varFF \in \HierLevS{1}$ et $\varFF^{\varCM}$. On
    suppose sans perte de généralité que $\varFF$ est de la forme
    $\exists x, \varFP$ où $\varFP\in\DeltaZL$. On sait donc qu'il existe un
    ensemble $a\in\varCM$ tel que $\varFP[a/x]$. Comme $\varFP[a/x]$ est
    $\DeltaZL$, elle est absolue entre ces deux classes transitives, et
    $\varCM \subseteq\varCM'$ donc $\varFP[a/x]$ permet de déduire
    $(\exists x, \varFP)^{\varCM'}$.
  \item supposons que $\varFF \in \HierLevP{1}$ et $\varFF^{\varCM'}$. On
    suppose sans perte de généralité que $\varFF$ est de la forme
    $\forall x, \varFP$ où $\varFP\in\DeltaZL$. On sait que donc pour tout
    ensemble $a\in \varCM'$, $\varFP[a/x]$. Ainsi, pour tout ensemble
    $a\in \varCM$, $\varFP[a/x]$ est vérifiée aussi, par inclusion. Donc
    $(\forall x, \varFP)^{\varCM}$.
  \end{itemize}

  Si $\varFF\in\HierLevD{1}$, alors $\varFF$ est $\HierLevS{1}$ et
  $\HierLevP{1}$, ce qui signifie que $\varFF^{\varCM}\iff \varFF^{\varCM'}$, donc
  que $\varCM\absolF \varCM'$.
\end{proof}

Donnons maintenant le théorème de réflexion, qui est un théorème majeur de la
théorie des ensembles. Celui-ci nous permet de montrer que pour toute hiérarchie
telle $\bV$ ou $\bL$ et toute proposition $\varFF$, il existe des ordinaux
$\varOb$ arbitrairement élevés tels que $\varFF$ est absolue entre la hiérarchie
et sa construction partielle à l'ordinal $\varOb$.

Ce théorème est plutôt un schéma de théorème~: pour chaque $\varFF$, il existe
un théorème correspondant. La démonstration, sans surprise, s'effectuera par
induction sur $\varFF$. Cependant, on remarque un point important~: pour une
formule sans quantificateur, $\varFF^\varCM = \varFF$. On n'a donc qu'à
traiter le cas des quantifications, et en utilisant la mise sous forme prénexe
des formules, déduire le résultat pour toute formule. De plus, en utilisant les
lois de De Morgan, on sait qu'on peut se contenter de prouver ce fait pour un
quantificateur (on choisit $\exists$) et de traiter le cas de la négation~:
celui-ci est direct, puisque si on sait que $\varFF \iff \varFP$ alors
$\lnot\varFF \iff \lnot \varFP$.

Pour montrer le résultat dans le cas de $\exists$, on a besoin d'un lemme sur
l'absoluité.

\begin{lemma}\label{lem.refl1}
  Soit une formule $\quantif_1\; x_1,\ldots,\quantif_k\;x_k, \varFP$ où les
  $\quantif_i$ sont des quantificateurs et $\varFP$ est sans quantificateurs,
  $\varCW$ une classe, $(\varCW_n)_{n\in \bN}$ une suite croissante d'ensembles
  inclus dans $\varCW$, et
  \[\varCW_\Ow \defeq \bigcup_{n \in \bN} \varCW_n\]
  Si pour tout $i \leq k$ et $n \in \bN$, on a
  \[\varCW_n \absol{\quantif_i\;x_i,\ldots,\quantif_k\;x_k,\varFP} \varCW\]
  alors on a aussi pour tout $i \leq k$
  \[\varCW_\Ow \absol{\quantif_i\; x_i,\ldots, \quantif_k\;x_k, \varFP}\varCW\]
\end{lemma}

\begin{proof}
  On raisonne par induction sur $k$, en ajoutant des négations ou des
  quantificateurs $\exists$~:
  \begin{itemize}
  \item si la formule est simplement $\varFP$, sans quantificateurs, alors comme
    $\varFP^\varCW = \varFP^{\varCW_\Ow}$, le résultat est automatiquement vrai.
  \item si la formule est $\lnot \varFF$ où $\varFF$ est une formule absolue
    entre $\varCW_\Ow$ et $\varCW$, comme on sait que $\varFC \iff \varFT$ est
    équivalente à $\lnot\varFC \iff \lnot \varFT$ pour toutes formules
    $\varFC,\varFT$, on en déduit à partir de l'hypothèse d'induction que
    $\lnot \varFF$ est absolue entre $\varCW_\Ow$ et $\varCW$ (et de même pour
    toutes les formules partielles obtenues à partir de $\varFF$ en en
    supprimant des quantificateurs en tête).
  \item si la formule est
    $\exists x, \quantif_1\;x_1,\ldots,\quantif_k\;x_k,\varFP$, alors en posant
    \[\varFF \defeq \quantif_1\;x_1,\ldots,\quantif_k\;x_k, \varFP\]
    on sait par hypothèse d'induction que chaque
    $\quantif_j\;x_j,\ldots,\quantif_k\;x_k,\varFP$ est absolue entre
    $\varCW_\Ow$ et $\varCW$.
    On souhaite maintenant montrer que $\exists x, \varFF$ est absolue
    entre $\varCW_\Ow$ et $\varCW$. Soient $p$ paramètres $\vec a \in X^p$,
    montrons
    \[\exists x\in \varCW_\Ow, \varFF^{\varCW_\Ow}(\vec a)
    \iff \exists x \in \varCW, \varFF^{\varCW}(vec a)\]
    par double implication~:
    \begin{itemize}
    \item supposons que $\exists x \in \varCW_\Ow,\varFF^{\varCW_\Ow}(\vec a)$.
      On trouve alors $b \in \varCW_\Ow$ vérifiant la formule
      $\varFF^{\varCW_\Ow}(b,\vec a)$. Par hypothèse d'induction, comme $\varFF$
      est absolue entre $\varCW_\Ow$ et $\varCW$, on en déduit que
      $\varFF^{\varCW}(b,\vec a)$, donc il existe bien
      $x \in \varCW$ tel que $\varFF^{\varCW}(\vec a)$.
    \item réciproquement, si
      $\exists x\in\varCW, \varFF^{\varCW}(\vec a)$, en prenant un tel $b$, on
      trouve $j$ tel que $\vec a \in \varCW_j$~: par absoluité entre $\varCW_j$
      et $\varCW$, on sait donc qu'il existe $b \in \varCW_j$, donc
      $b \in \varCW_\Ow$, tel que $\varFF^{\varCW}(b,\vec a)$ donc, par absoluité
      de $\varFF$ entre $\varCW_\Ow$ et $\varCW$, on en déduit
      que $\exists x \in \varCW_\Ow, \varFF^{\varCW}(vec a)$.
    \end{itemize}
    On en déduit donc que $\exists x, \varFF$ est absolue entre $\varCW_\Ow$
    et $\varCW$.\qedhere
  \end{itemize}
\end{proof}

\noindent
On peut maintenant montrer le théorème à proprement parler.

\begin{theorem}[Réflexion]\label{thm.reflexion}
  Soit une fonction $\varCW : \Ord \to \mathbb V$ vérifiant~:
  \begin{itemize}
  \item pour tous $\varOa, \varOb \in \Ord$,
    $\varOa \leq \varOb \implies \varCW_\varOa \subseteq \varCW_\varOb$.
  \item pour tout $\varOl \in \Ord$ limite,
    $\displaystyle \varCW_\varOl = \bigcup_{\varOb < \varOl} \varCW_\varOb$
  \end{itemize}
  et $\displaystyle \varCW = \bigcup_{\varOa \in \Ord} \varCW_\varOa$.
  alors pour toute proposition $\varFF\in\Formula(\SignP{\ZF})$ et tout
  ordinal $\varOa \in \Ord$, il existe un ordinal $\varOb > \varOa$ tel que
  $\varCW_\varOb\absolF \varCW$. On peut de plus prendre $\varOb$ limite.
\end{theorem}

\begin{proof}
  On montre le résultat uniquement sur les formules sous forme prénexe, mais
  comme toute formule est équivalente à une formule sous forme prénexe, cela
  suffit à montrer le résultat souhaité. On procède donc par récurrence sur
  le nombre de quantificateurs de la formule, en ajoutant des $\exists$ et des
  négations~:
  \begin{itemize}
  \item comme dans le lemme précédent, le cas sans quantificateur se traite
    directement puisque la relativisation ne change pas la formule.
  \item dans le cas d'une négation, comme $\varFF \iff \varFP$ et
    $\lnot\varFF\iff\lnot\varFP$ sont équivalentes, on a aussi le résultat
    directement.
  \item on suppose donc que la formule pour laquelle on souhaite prouver notre
    énoncé est $\exists x, \varFF(x,x_1,\ldots,x_n)$. Par hypothèse
    d'induction, on sait déjà que pour tout $\varOa \in \Ord$, on peut trouver
    $\varOb > \varOa$ limite, tel que $\varCW_\varOb \absolF \varCW$.

    On définit maintenant une fonction (entre classes) $F$, à $n$ arguments,
    donnée par le fait que $F(x_1,\ldots,x_n)$ est l'ensemble des
    $x \in \varCW$ tels que $\varFF(x,x_1,\ldots,x_n)$ et dont le
    $\varCW$-rang (au sens de plus petit $\varOa$ tel que l'élément est dans
    $\varCW_\varOa$) est minimal. Ainsi, quitte à choisir uniquement les éléments
    de $\varCW$-rang minimal, on a l'équivalence suivante, pour toute formule
    $\varFP$~:
    \[\exists x \in \varCW, \varFP(x,x_1,\ldots,x_n) \iff
    \exists x \in F(x_1,\ldots,x_n), \varFP(x,x_1,\ldots,x_n)\]

    Soit $\varOa \in \Ord$. On décide maintenant de définir une suite
    $(\varOb_k)_{k \in \bN}$ d'ordinaux limites supérieurs à $\varOa$, par
    récursion d'ordre 2~:
    \begin{itemize}
    \item on choisit pour $\varOb_0$ le premier $\varOb > \varOa$ tel que
      $\varCW_\varOb \absolF \varCW$.
    \item pour un rang $k$ fixé, on définit $\varOb_{2k + 1}$ comme le plus
      petit ordinal strictement supérieur à $\varOb_{2k}$ tel que
      $\varCW_{\varOb_{2k + 1}}$ contienne $F(a_1,\ldots,a_n)$ pour tous
      $a_1,\ldots,a_n \in \varCW_{\varOb_{2k}}$ (il nous suffit de considérer
      l'union de tous ces $F(a_1,\ldots,a_n)$ pour obtenir un ensemble, dont les
      $\varCW$-rangs sont donc majorés, et de prendre un ordinal supérieur
      à cette majoration).
    \item on définit ensuite $\varOb_{2k + 2}$ comme le plus petit ordinal
      supérieur strictement à $\varOb_{2k + 1}$ tel que
      $\varCW_{\varOb_{2k + 2}}\absolF \varCW$.
    \end{itemize}
    On définit alors $\varOb = \sup \varOb_k$~: c'est un ordinal limite car
    les $\varOb_k$ sont des ordinaux limites. D'après le \cref{lem.refl1}, et
    comme les $\varCW_{\varOb_k}$ forment une suite croissante et incluse dans
    $\varCW$ par hypothèse, on sait donc que $\varFF$ est absolue entre
    $\varCW_\varOb$ et $\varCW$ (puisque $\varOb$ est le sup des $\varOb_{2k}$).
    On montre maintenant que $\exists x, \varFF$ est absolue entre
    $\varCW_\varOb$ et $\varCW$. Soient $\vec a \in (\varCW_\varOb)^n$, ce que
    l'on peut réécrire par $vec a \in (\varCW_{2k})^n$ pour un certain
    $k \in \bN$ (que l'on fixe désormais). Montrons par double implication que
    \[\exists x\in \varCW_\varOb, \varFF^{\varCW_\varOb}(\vec a) \iff
    \exists x \in \varCW, \varFF^{\varCW}(\vec a)\]
    \begin{itemize}
    \item si on trouve $b \in \varCW_\varOb$ vérifiant l'énoncé, alors en
      utilisant le fait que $\varFF$ est absolue entre $\varCW_\varOb$ et
      $\varCW$ et le fait que $\varCW_\varOb \subseteq \varCW$, on obtient le
      résultat.
    \item réciproquement, si l'on trouve $b \in \varCW$ tel que
      $\varFF^\varCW(b,\vec a)$ alors on peut aussi le trouver de $\varCW$-rang
      minimal. Cela signifie donc que $b \in F(\vec a)$, donc on sait que
      $b \in \varCW_{\varOb_{2k + 1}}$, donc $b \in \varCW_\varOb$. On en déduit,
      comme $\varFF$ est absolue entre $\varCW_\varOb$ et $\varCW$, qu'il existe
      $b\in \varCW_\varOb$ tel que $\varFF^{\varCW_\varOb}(b,\vec a)$, d'où le
      résultat.\qedhere
    \end{itemize}
  \end{itemize}
\end{proof}

\begin{remark}
  La preuve a été légèrement allégée pour être plus digeste~: pour vérifier les
  hypothèses du \cref{lem.refl1}, il nous faut établir comme énoncé à montrer
  par induction que le $\varOb$ que l'on trouve à chaque fois est vrai non
  seulement pour $\varFF$, mais aussi pour $\varFF$ auquel on enlève une
  partie des quantificateurs de tête. Remarquons aussi que, comme le
  \cref{lem.refl1} est lui-même un schéma de lemme, et que l'on veut prouver
  un schéma de théorèmes, il est bon de vérifier un fait~: pour une formule
  $\varFF$ donnée, il nous suffit d'utiliser notre schéma de lemmes un nombre
  fini de fois pour atteindre la conclusion.
\end{remark}

\begin{remark}
  Comme on peut toujours représenter un ensemble fini de propositions $F$ par
  $\bigwedge F$, le théorème de réflexion s'étend directement à un ensemble fini
  de propositions.
\end{remark}

Une première conséquence est un résultat assez peu surprenant, mais qu'il est
intéressant d'avoir.

\begin{proposition}
  La théorie $\ZFC$ n'est pas finiment axiomatisable.
\end{proposition}

\begin{proof}
  Si $\ZFC$ était finiment axiomatisable, alors on trouverait un certain
  $\bV_\varOa$ tel que $\ZFC$ est absolu pour $\bV_\varOa$, nous
  donnant donc un modèle ensembliste de $\ZFC$.
\end{proof}

\section{L'univers constructible de Gödel}

Nous avons introduit cet univers dans la \cref{def.L}, mais ne l'avons pas
encore étudié. Nous le verrons, c'est un univers ensembliste se comportant
particulièrement bien~: l'axiome du choix y est vrai dans une version
particulièrement forte et l'hypothèse du continu généralisée y est vraie
(d'autres propriétés qui dépassent le cadre de ce cours y sont vraies aussi).
Nous allons donc étudier plus en profondeur l'univers $\bL$.

\subsection{Premiers pas}

Tout d'abord, intéressons-nous à l'opération $\powerdef$. Celle-ci ressemble
à l'opération $\powerset$, mais on peut remarquer que son action sur le cardinal
n'est pas du tout la même.

\begin{proposition}
  Les assertions suivantes sont vérifiées~:
  \begin{itemize}
  \item si $X$ est un ensemble fini, alors $\powerdef(X) = \powerset(X)$.
  \item si $X$ est un ensemble infini, alors $\Card(\powerdef(X)) = \Card(X)$.
  \end{itemize}
\end{proposition}

\begin{proof}
  Prouvons ces deux assertions~:
  \begin{itemize}
  \item La première inclusion, $\powerdef(X)\subseteq\powerset(X)$, est directe.
    Supposons que $X$ est fini et soit $Y\subseteq X$. Décomposons
    $Y = \{y_1,\ldots,y_p\}$. On peut alors définir la proposition
    \[\varFF(x) \defeq \bigvee_{i = 1}^p (x = y_i)\]
    dont il est évident que $x \in Y \iff X \satisf \varFF(x)$. Ainsi
    $Y\in \powerdef(X)$, d'où l'égalité.
  \item Supposons que $X$ est infini. Pour énumérer $\powerdef(X)$, il nous
    suffit d'énumérer $\Formula(\SignP{\ZF})$, qui est dénombrable, et
    l'ensemble des parties finies de $X$~: $\Card(X^{<\Ow}) = \Card(X)$ car
    $X$ est au moins dénombrable. Ainsi il y a
    $\Card(\Ow \times X^{<\Ow}) = \Card(X)$ tuples $(\varFF,\vec a)$ définissant
    une partie définissable de $X$, donc $\Card(\powerdef(X))\leq |X|$. De plus,
    on peut injecter $X$ dans $\powerdef(X)$ par $x\mapsto \{x\}$, dont la
    représentation est la formule
    \[\varFP(y) \defeq x = y \]
    donc $\Card(\powerdef(X)) = \Card(X)$.\qedhere
  \end{itemize}
\end{proof}

On montre aussi, de façon analogue au travail effectué pour $\bV$, que
$(\bLa)_{\varOa \in \Ord}$ est une hiérarchie cumulative et transitive.

\begin{proposition}
  Pour tout $\varOa,\varOb \in \Ord$, si $\varOa \leq \varOb$ alors
  $\bLa \subseteq \bL_\varOb$. De plus chaque $\bL_\varOa$ est un ensemble
  transitif.
\end{proposition}

\begin{proof}
  On prouve d'abord que chaque $\bLa$ est un ensemble transitif, par induction
  sur $\varOa$~:
  \begin{itemize}
  \item $\bL_0 = \vide$ est transitif.
  \item supposons que $\bLa$ est transitif. Soit alors
    $x \in \bL_{\varOa + 1}$. Cela signifie que
    $x\subseteq \bL_{\varOa + 1}$, donc $y \in \bL_\varOa$ pour tout
    $y \in x$. Mais alors on peut définir $y$ par la formule
    \[\varFF_y(x) \defeq x \in y\]
    à paramètre ($y$) dans $\bLa$. Ainsi $y \in \bL_{\varOa + 1}$, donc
    $x \subseteq \bL_{\varOa + 1}$.
  \item pour le cas limite, l'union d'ensembles transitifs est un ensemble
    transitif.
  \end{itemize}

  On prouve maintenant que la hiérarchie est cumulative, par induction sur
  la différence entre $\varOb$ et $\varOa$~:
  \begin{itemize}
  \item $\bLa \subseteq \bL_\varOa$
  \item on montre que $\bL_{\varOb}\subseteq \bL_{\varOb + 1}$, ce qui suffit à
    montrer que
    $\bLa \subseteq \bL_\varOb \implies \bL_\varOa \subseteq \bL_{\varOb + 1}$.
    Si $x \in \bL_{\varOb}$, alors on peut définir $x$ par
    \[\varFF_x(a) \defeq a \in x\]
    dans $\bL_{\varOb + 1}$, donc $x \in \bL_{\varOb + 1}$.
  \item le cas où la différence est un ordinal limite est direct.
  \end{itemize}
  Ainsi, par induction transfinie, $\bLa$ est une hiérarchie
  cumulative.
\end{proof}

Il est donc clair que $\powerdef(X)\neq \powerset(X)$ lorsque $X$ est infini.
On peut donc montrer que $\bL$ ne coïncide avec $\bV$ que jusqu'à $\Ow$.

\begin{proposition}
  Pour tout $\varOa \leq \Ow$, $\bLa = \bVa$. Cependant, pour tout
  $\varOa > \Ow$, on a seulement l'inclusion $\bLa \subseteq \bVa$. Lorsque
  $\varOa > \Ow^2$, si $\Card(\varOa) \neq \Cbe_\varOa$ alors
  $\bLa \subsetneq \bVa$.
\end{proposition}

\begin{proof}
  On procède par induction transfinie~:
  \begin{itemize}
  \item $\bL_0 = \bV_0$
  \item pour tout $n \in \Ow$, si $\bL_n = \bV_n$, alors comme
    les deux ensembles sont finis~:
    \[\bL_{n+1} = \powerdef(\bL_n)
      = \powerset(\bL_n)
      = \powerset(\bV_n)
      = \bV_{n+1}\]
  \item comme $\forall n < \Ow, \bL_n = \bV_n$, on en déduit que
    $\bL_\Ow = \bigcup \bL_n = \bV_\Ow$.
  \item comme $\bL_\Ow$ est infini,
    $\bL_{\Ow+1} \subsetneq \powerset(\bL_\Ow)$, donc
    $\bL_{\Ow+1} \subsetneq\bV_{\Ow+1}$.
  \item comme on sait que $\Card(\powerdef(X)) = \Card(X)$ pour $X$ infini, on
    en déduit que $\Card(\bLa) = \Card(\varOa)$ pour $\varOa > \Ow$. En
    comparaison, $\Card(\bVa) = \Cbe_\varOa$ à partir de $\varOa^2$, donc
    l'égalité des deux ensembles implique que $\Card(\varOa) = \Cbe_\varOa$, d'où
    l'inclusion stricte par contraposée.
  \end{itemize}
\end{proof}

\noindent
Cependant, $\bL$ contient bien tous les ordinaux.

\begin{proposition}
  Pour tout $\varOa \in \Ord$, $\Ord \cap \bLa = \varOa$.
\end{proposition}

\begin{proof}
  On raisonne par induction transfinie~:
  \begin{itemize}
  \item le cas de $\vide$ est direct.
  \item on suppose que $\Ord \cap \bL_\varOa = \varOa$, montrons que
    $\Ord\cap \bL_{\varOa + 1} = \varOa + 1$ par double inclusion.

    Soit $\varOb \in\Ord\cap \bL_{\varOa + 1}$~: on sait donc que $\varOb$ est une
    partie définissable de $\bLa$~: en particulier, c'est un ordinal appartenant
    à $\powerset(\varOa)$ puisque les ordinaux dans $\bLa$ sont exactement les
    éléments de $\varOa$. On en déduit que $\varOb < \varOa + 1$.

    Soit $\varOb < \varOa + 1$. Si $\varOb < \varOa$, alors $\varOb \in \bLa$
    par hypothèse d'induction, donc $\varOb \in \bL_{\varOa + 1}$ par croissance
    de la hiérarchie $\bL$. Si $\varOb = \varOa$, alors on considère la formule
    \[\varFP(x)\defeq \Ord(x)\]
    On sait que $\Ord\cap \bLa = \varOa$, donc
    $(\bLa \satisf \varFP(x)) \iff x \in \varOa$, d'où le fait que $\varOa$ est
    une partie définissable de $\bLa$, donc un élément de $\bL_{\varOa + 1}$.
  \item dans le cas où $\varOl$ est un ordinal limite, on suppose que pour tout
    $\varOb < \varOl$, $\Ord \cap \bL_\varOb = \varOb$, mais alors l'équation
    passe à la borne supérieure pour donner $\Ord\cap \bL_\varOl = \varOl$ par
    continuité de $\bL$.
  \end{itemize}
  Ainsi pour tout $\varOa \in \Ord, \Ord \cap \bLa = \varOa$.
\end{proof}

\begin{corollary}
  Pour tout $\varOa \in \Ord, \varOa \in \bL_{\varOa + 1}$, et
  $\Ord \subseteq \bL$.
\end{corollary}

Il y a donc toujours plus d'éléments dans $\bVa$ que dans $\bLa$. Cependant, on
pourrait imaginer qu'en prenant par exemple $\varOa' = \Cbe_\varOa$, on puisse
rattraper $\bVa$ en considérant $\bL_{\varOa'}$. Le fait que l'on puisse toujours
rattraper la hiérarchie $(\bVa)_{\varOa\in\Ord}$ avec $\bL$ est une assertion
indépendante de $\ZFC$, qu'on introduit sous la forme de l'axiome suivant~:

\begin{axiom}[V = L]\label{ax.VeqL}
  On définit l'axiome $\VeqL$ par la formule suivante~:
  \[\VeqL \defeq \forall x, \exists \varOa \in \Ord, x \in \bLa\]
\end{axiom}

Nous étudions maintenant les propriétés de $\ZFC+\VeqL$. Pour cela, cherchons
d'abord à situer dans la hiérarchie de Lévy les éléments importants de l'univers
constructible.

\begin{property}
  Pour une relation fonctionnelle $H$ telle qu'introduite dans la
  \cref{subsec.ind.transf}, si $H$ est $\DeltaZL$ alors pour tout ordinal
  $\varOa$ et toute fonction $f : \varOa \to \bV$, le prédicat
  $\ind_H(f)$ est $\DeltaZL$. La relation fonctionnelle définie par récursion
  transfinie par $H$ est une relation $\HierLevD{1}$.
\end{property}

\begin{proof}
  En effet, dire que $f : \varOa \to \bV$ est $H$-inductive signifie que~:
  \begin{itemize}
  \item $\forall \varOb < \varOa, f\restr{\varOb}\subseteq \dom(H)$
  \item $\forall \varOb < \varOa, f(\varOb) = H(f\restr{\varOb})$
  \end{itemize}
  Il nous suffit donc de montrer~:
  \begin{itemize}
  \item qu'être le domaine de $f$ est un prédicat $\DeltaZL$, et on peut
    exprimer \og $X$ est le domaine de $f$\fg par
    \[\forall (x,y)\in f, x\in X \land \forall x \in X, \exists (z,y)\in f,
    x = z\]
    donc $\dom(f)$ est défini par un prédicat $\DeltaZL$.
  \item qu'être une restriction de $f$ est un prédicat $\DeltaZL$, ce qui est
    direct en voyant que
    \[g = f\restr{X} \iff \dom(g)\subseteq\dom(f)\land
    \forall x \in \dom(g), g(x)=f(x)\]
  \end{itemize}
  On sait donc que, pour une fonction $f : \varOa \to \bV$, savoir si
  $f$ est $H$-inductive est $\DeltaZL$.

  \'Etant donné $H$, on considère la relation fonctionnelle
  \[F(\varOa,y) \defeq \exists f : \Ss \varOa \to \bV,
  \ind_H(f) \land (\varOa,y) \in f\]
  qui définit l'unique fonction $H$-inductive $\Ord \to \bV$. Mais, puisqu'il
  y a unicité d'une telle fonction $H$-inductive pour un ordinal $\varOa$
  donné, on peut tout aussi bien définir notre relation fonctionnelle par
  \[F(\varOa,y) \defeq \forall f : \Ss \varOa \to \bV,
  \ind_H(f) \implies (\varOa,y) \in f\]
  donc la fonction $\Ord \to \bV$ définie par $H$ est, elle, $\HierLevD{1}$.
\end{proof}

Ainsi, exprimer le prédicat d'être $\bL_\varOa$ (pour un ordinal $\varOa$),
est à un niveau comparable à celui du prédicat utilisé pour sa définition, qui
est clairement la complexité de $\powerdef$.

C'est précisément ici que l'on voit que $\bL$ se comporte beaucoup mieux
que $\bV$~: là où $\powerset$ est une fonction $\HierLevP{1}$, $\powerdef$ n'est
que $\HierLevD{1}$, qui est une formule absolue.

\begin{property}\label{prop.Form.Delta}
  Le prédicat $Y = \Formula(\SignP{\ZF})$ est $\DeltaZL$.
\end{property}

\begin{proof}
  La définition est par induction, et la relation permettant de construire
  de nouvelles formules à partir d'anciennes est clairement $\DeltaZL$.
\end{proof}

\begin{property}
  Le prédicat $Y = \powerdef(X)$ est $\DeltaZL$.
\end{property}

\begin{proof}
  On sait grâce à la \cref{prop.Form.Delta} que l'on peut quantifier
  sur $\Formula(\Sign{\ZF})$ sans souci. La définition de $X\satisf \varFF$ est
  une induction (non transfinie~!) sur $\Formula(\SignP{\ZF})$. La
  quantification sur les éléments de $X$ est, elle aussi, $\Delta_0$. Ainsi
  \begin{multline*}
    \exists \varFF(\vec x,y)\in\Formula_{n+1}(\SignP{\ZF}),
    \exists \vec a\in X^n, \forall b\in X, \\
    b\in Y \iff X\satisf \varFF(\vec a,b)
  \end{multline*}
  est bien $\DeltaZL$, donc $Y\in \powerdef(X)$ est un prédicat $\DeltaZL$.
  Maintenant, dire que $Y = \powerdef(X)$ signifie
  \begin{multline*}
    (\forall y \in Y, y \in \powerdef(X)) \land
    \forall \varFF\in\Formula_{n+1}(\SignP{\ZF}),
    \forall \vec a \in X^n, \exists y \in Y, \forall b \in X,\\
    b \in y \iff X\satisf \varFF(\vec a,b)
  \end{multline*}
  donc $Y = \powerdef(X)$ est un prédicat $\DeltaZL$.
\end{proof}

\begin{corollary}
  Pour tout ordinal $\varOa$ et toute classe transitive $\varCM$, le
  prédicat $x\in \bL_{\varOa}$ est absolu sur $\varCM$ en tant que prédicat
  $\HierLevD{1}$.
\end{corollary}

\noindent
On peut donc placer l'axiome de constructibilité dans la hiérarchie de Lévy.

\begin{proposition}
  $\VeqL \in \HierLevP{2}$.
\end{proposition}

\begin{proof}
  On peut énoncer $\VeqL$ par $\forall x, \exists \varOa \in \Ord, x\in \bLa$
  qui est $\HierLevP{2}$ avec les résultats précédents.
\end{proof}

Une conséquence intéressante est que $\VeqL$ est un axiome vérifié dans $\bL$.
Cela parait assez évident, mais il faut rappeler qu'\latinexpr{a priori}, on
pourrait obtenir plus de parties en regardant $\powerdef$ dans $\bV$.
L'avantage justement ici est que toutes nos fonctions sont absolues pour $\bL$. 

\begin{theorem}
  $\ZFC\satisf (\VeqL)^{\bL}$
\end{theorem}

\begin{proof}
  On sait que le prédicat $x\in \bL_{\varOa}$ est vrai dans $\bL$ si et seulement
  s'il est vrai dans $\bV$. Par définition, on a
  \[(\VeqL)^{\bL}= \forall x\in\bL,\exists \varOa \in \Ord\cap \bL,x\in \bLa\]
  mais on a vu que $\Ord\subseteq \bL$, donc il nous reste à prouver que
  $\forall x \in \bL, \exists \varOa \in \Ord, x \in \bLa$ ce qui est vrai dans
  $\bV$ par définition, et donc dans $\bL$ par absoluité.
\end{proof}

On montre aussi que $\bL$ est un modèle de $\ZF$. On donne d'abord un
lemme, similaire à \cref{lem.axZFC}, mais utilisant des conditions plus fortes.

\begin{lemma}\label{lem.axZFC2}
  Soit $\varCM$ une classe transitive vérifiant l'axiome de compréhension et
  telle que pour tout $x \subseteq \varCM$, si $x$ est un ensemble, alors il
  existe $y \in \varCM$ tel que $x \subseteq y$. Dans ce cas, $\varCM$ est un
  modèle de $\ZF$.
\end{lemma}

\begin{proof}
  On traite les différents axiomes~:
  \begin{itemize}
  \item axiome d'extensionalité~: $\varCM$ est une classe transitive.
  \item axiome de fondation~: vrai puisque $\varCM$ est une partie de $\bV$.
  \item axiome de la paire~: on voit que si $x, y \in \varCM$ alors
    $\{x,y\} \subseteq \varCM$, puis on applique l'axiome de compréhension
    à l'ensemble contenant $\{x,y\}$ et appartenant à $\varCM$ qu'on trouve par
    hypothèse sur $\varCM$.
  \item axiome de la réunion~: si $x \in \varCM$ alors par transitivité de
    $\varCM$, $\bigcup x \subseteq \varCM$, d'où le résultat par compréhension
    et hypothèse sur $\varCM$.
  \item axiome de compréhension~: vrai par hypothèse.
  \item axiome de remplacement~: supposons que $f$ est une fonction dont le
    domaine est $x \in \varCM$ et l'image $y \subseteq \varCM$. Alors par
    hypothèse on trouve $z \in \varCM$ contenant $y$, et on utilise l'axiome de
    compréhension sur $z$ avec la formule
    \[\varFF(b) \defeq \exists a \in x, f(x) = b\]
    pour en déduire que $y \in \varCM$.
  \item axiome de l'ensemble des parties~: soit $x \in \varCM$. On
    considère
    \[\powerset_\varCM(x) \defeq \compr{y \in \varCM}{y \subseteq x}\]
    par hypothèse sur $\varCM$, on trouve $z \in \varCM$ tel que
    $\powerset_\varCM(x) \subseteq z$, et par compréhension on peut restreindre
    $z$ aux éléments contenus dans $x$, nous redonnant $\powerset_\varCM(x)$.
    Ainsi pour tout
    $y \in \varCM, y \in \powerset_{\varCM}(x) \iff y \subseteq x$.
  \item axiome de l'infini~: comme $\varCM$ est clos par parties finies grâce à
    l'axiome de la paire et de la réunion, et puisque $\vide \in \varCM$ par
    compréhension, on en déduit que pour tout $n \in\Ow,\bV_n\subseteq\varCM$,
    donc $\bV_\Ow \subseteq \varCM$. On applique alors notre hypothèse pour
    trouver $z$ contenant $\bV_\omega$ et donc $\Ow$, puis on utilise la
    compréhension.
  \end{itemize}
  On a ainsi montré notre critère pour avoir un modèle de $\ZF$.
\end{proof}

\noindent
On peut ainsi prouver que $\bL$ est un modèle (classe) de $\ZF$.

\begin{theorem}
  $\ZF\satisf (\ZF)^{\bL}$.
\end{theorem}

\begin{proof}
  On vérifie les hypothèses du \cref{lem.axZFC2}. On sait déjà que $\bL$ est
  transitif. Pour la dernière hypothèse, si $x$ est un ensemble inclus dans
  $\bL$, alors pour tout $y \in x$ on trouve $\varOa_y$ tel que
  $y \in \bL_{\varOa_y}$. On sait alors que pour
  \[\varOa \defeq \sup_{y \in x} \varOa_y\]
  on a $x \subseteq \bL_\varOa \in \bL$.

  On veut maintenant prouver l'axiome de compréhension. Soit donc
  $\varFP^{\bL}(\vec z,x,y)$ une formule sur $\SignP{\ZF}$, et $\vec a,b$ des
  paramètres dans $\bL$. Quitte à prendre une borne supérieure, soit $\varOa$
  tel que $\vec a,b \in \bLa$. Par le \cref{thm.reflexion}, on trouve
  $\varOb > \varOa$ tel que $\varFP^{\bL}(\vec z,x,y)$ est absolue entre
  $\bL_\varOb$ et $\bL$. On a alors
  \[b\defeq \compr{c \in \bL_\varOb}{c \in b \land
  \varFP^{\bL}(\vec a,b,c)} \in \bL_{\varOb + 1}\]
  et un élément de $\bL$ vérifie $\varFP^{\bL}$ si et seulement s'il le vérifie
  dans $\bL_{\varOb}$ (car pour vérifier $\varFP^{\bL}$, il lui faut au moins être
  un élément de $b$, donc de $\bL_\varOb$), si et seulement s'il appartient à
  $b$. On a donc vérifié l'axiome de compréhension.

  Ainsi $\bL \satisf \ZF$.
\end{proof}

Il en découle donc que $\VeqL$ est relativement cohérent~: on peut ajouter cet
axiome sans risque puisque si ajouter $\VeqL$ mène à une incohérence, alors
cette incohérence était déjà dans $\ZF$.

\begin{corollary}
  Si $\Coher(\ZF)$ alors $\Coher(\ZF + \VeqL)$.
\end{corollary}

\begin{remark}
  On verra qu'en fait $\VeqL$ implique une version forte de l'axiome du choix,
  donc on a aussi une preuve de $\Coher(\ZF)\implies \Coher(\ZFC)$ puisque nos
  constructions sur $\bL$ n'ont pas nécessité l'axiome du choix.
\end{remark}

\subsection{Les propriétés de l'univers constructible}

On montre dans cette sous-section que l'univers $\bL$ vérifie des résultats
particulièrement forts~: l'axiome du choix global et l'hypothèse du continu
généralisée. On suppose donc l'axiome $\VeqL$ en plus des autres axiomes de
$\ZFC$.

Présentons donc l'axiome du choix global. Plutôt qu'un simple axiome, celui-ci
a besoin d'autres données pour être exprimé. Nous considérons que nous sommes
munis d'un symbole $\sCh$ de fonction unaire.

\begin{axiom}[Axiome du choix global]
  L'axiome du choix global, pour le symbole de fonction $\sCh$, s'exprime par
  \[\forall x, x\neq \vide \implies \sCh(x)\in x\]
\end{axiom}

Dans notre cas, puisqu'on cherche à représenter des fonctions par des prédicats
binaires fonctionnels, on peut à la place considérer une formule $\varFF(x,y)$
fonctionnelle. L'axiome du choix global pour $\varFF$ (supposée fonctionnelle)
est alors
\[\forall x, x \neq \vide \implies (\forall y, \varFF(x,y)\implies y\in x)\]
On cherche donc à construire une telle formule $\varFF$.

En fait, on voit avec les constructions précédentes à propos de l'axiome du
choix qu'il nous suffit de pouvoir bien ordonner l'univers, c'est-à-dire
construire un prédicat $\leL$ vérifiant les axiomes suivants~:
\begin{itemize}
\item $\forall x, \lnot(x \leL x)$
\item $\forall x,y,z, (x \leL y \land y \leL z) \implies x \leL z$
\item $\forall x,y, x\neq y \implies x \leL y \lor y \leL x$
\item $\forall x, x\neq \vide \implies
  \exists y, y \in x \land (\forall z \in x, \lnot (z \leL y))$
\end{itemize}
La fonction est alors simplement
\[\varFF(x,y) \defeq y \in x \land \forall z \in x, \lnot(z \leL y)\]
qui est évidemment fonctionnelle totale pour $x\neq \vide$ puiqu'il existe un
tel $y$ par le dernier point, et les autres axiomes assurent que ce $y$ est
unique.

Il nous faut donc construire $\leL$. Pour cela, remarquons deux choses~:
\begin{itemize}
\item $\Formula(\SignP{\ZF})$ peut être bien ordonné.
\item si $\bLa$ peut être bien ordonné, alors l'ensemble des suites finies
  ${\bLa}^{<\omega}$ à valeurs dans $\bLa$ peut l'être aussi.
\end{itemize}
En utilisant ces deux faits, comme un élément de $\bL_{\varOa + 1}$
peut être associé à un tuple $(\varFF,\vec x)$ où
$\varFF\in\Formula_{n+1}(\SignP{\ZF})$ et $\vec x \in(\bLa)^n$, on peut bien
ordonner $\bL_{\alpha + 1}$. Remarquons cependant que plusieurs tuples peuvent
correspondre à une seule partie définissable~: il faut donc prendre à chaque
fois le tuple le plus petit possible.

Enfin, la question se pose du recollement des ordres~: si on a défini un ordre
sur $\bLa$ il faut que celui défini sur $\bL_{\varOa + 1}$ contienne l'ordre
défini juste avant. Cela nous mène à l'énoncé que l'on veut donner.

\begin{theorem}
  Il existe une fonction $\mathbin{\leL} : \Ord \to \bV$ telle que que
  $\mathbin{\leL}(\varOa)$ est un bon ordre sur $\bL_\varOa$ pour tout
  $\varOa\in \Ord$ et telle que pour tous $\varOa < \varOb$,
  $\mathbin{\leL}(\alpha)$ est un segment initial de $\mathbin{\leL}(\beta)$.
\end{theorem}

\begin{proof}
  Notre démonstration est une induction transfinie~:
  \begin{itemize}
  \item pour $\vide$, le résultat est direct.
  \item supposons que $\mathbin{\leL}(\varOa)$ existe. On utilise ce qu'on a dit
    précédemment~: on se fixe un bon ordre $<$ sur
    $\Formula(\SignP{\ZF})\times {\bLa}^{<\omega}$, soit
    $\varOb$ le type de cet ordre. On peut associer à chaque partie
    $X\in\bL_{\varOa + 1}\setminus\bL_\varOa$ un élément de $\varOb$ comme étant le
    plus petit tuple $(\varFF,\vec a)$ tel que
    $x\in X \iff \bLa\satisf \varFF(\vec a,x)$.
    Cette association est clairement injective, vu que la partie définie par
    un tuple est unique. On a donc une injection de
    $\bL_{\varOa + 1}\setminus \bL_\varOa$ vers un ordinal, donc un bon ordre sur
    $\bL_{\varOa + 1} \setminus \bL_\varOa$. En prenant simplement la somme
    ordinale avec l'ordre défini sur $\bLa$, on trouve ainsi un bon ordre sur
    $\bL_{\varOa + 1}$ tel que l'ordre sur $\bLa$ en est un segment initial.
  \item si $\mathbin{\leL}(\varOb)$ a été défini pour tout ordinal
    $\varOb < \varOl$, avec $\varOl$ un ordinal limite, alors on peut définir
    $\mathbin{\leL}(\varOl)$ comme l'union des $\mathbin{\leL}(\varOb)$. Comme
    tous les bons ordres sont compatibles, l'union est encore un bon ordre.
  \end{itemize}
  On a donc construit notre fonction $\mathbin{\leL} : \Ord \to \bV$ telle que
  $\mathbin{\leL}(\varOa)$ est un bon ordre sur $\bLa$ pour tout
  $\varOa\in\Ord$.
\end{proof}

\noindent
On en déduit l'axiome du choix global.

\begin{corollary}
  Il existe une relation $\leL$ sur $\bL$ qui est un bon ordre au sens
  défini plus tôt, donc $\ZF\satisf (\AxC)^{\bL}$.
\end{corollary}

On veut maintenant montrer que $\bL$ vérifie $\HCG$.
Pour cela, on va montrer d'abord que, sous l'hypothèse $\VeqL$, la hiérarchie
$\bH$ coïncide avec $\bL$ sur les cardinaux infinis.

On remarque que, comme $\bH_\varCCl$ et $\bL_\varCCl$ peuvent s'exprimer comme
l'union des $\bH_\varCk$ (respectivement $\bL_\varCk$) pour $\varCk < \varCCl$
lorsque $\varCCl$ est un cardinal limite, il nous suffit de nous occuper du cas
successeur.

L'inclusion qui nous demandera du travail est $\bH_\varCk \subseteq\bL_\varCk$,
nous allons donc l'étudier à part. Pour traiter cette inclusion, on commence
par donner le lemme de condensation.

\begin{lemma}[Condensation]
  Soit $X$ un ensemble transitif qui est un sous-modèle élémentaire d'un
  ensemble $\bLa$ (pour $\varOa \in \Ord$), c'est-à-dire tel que pour toute
  proposition $\varFF\in\Prop(\SignP{\ZF})$,
  $X\satisf \varFF\iff \bLa\satisf\varFF$. Alors $X = \bL_{\varOb}$ pour un
  certain ordinal $\varOb \leq \varOa$.
\end{lemma}

\begin{proof}
  On veut prouver que $X = \bL_{\varOb}$ pour $\varOb \defeq X \cap \Ord$. On
  remarque que, puisqu'il est possible d'énoncer que $\Ord$ est (ou non) une
  classe bornée, $\varOb$ est un successeur si et seulement si $\varOa$ est un
  successeur (et de même pour les limites). On va procéder par double
  inclusion~:
  \begin{itemize}
  \item[\fbox{$\subseteq$}] On sait que $X\satisf\VeqL$, d'où
    \[(\VeqL)^X \iff\forall x\in X,\exists\varOc\in\Ord\cap X, x\in\bL_{\varOc}\]
    et, comme on a pris $\varOb \defeq X \cap \Ord$, on sait on peut en
    déduire que $\forall x \in X, \exists \varOc < \varOb, x \in \bL_{\varOc}$,
    donc $X \subseteq \bL_{\varOb}$.
  \item[\fbox{$\supseteq$}] On commence par utiliser le fait que la relation
    fonctionnelle $(\varOc,\bL_{\varOc})$ est absolue entre les ensembles
    transitifs (remarquons que pour construire cette relation, il nous suffit
    d'utiliser les axiomes de $\ZF - \AxPow$, et que les $\bL_{\varOc}$ sont des
    modèles de cette théorie). Par cette absoluité, et l'axiome de remplacement,
    on en déduit que $\bL_{\varOc} \in X$ pour tout $\varOc < \varOb$. Mais,
    puisque $X$ est transitif, on a $\bL_{\varOc} \subseteq X$ donc, par union
    pour $\varOa < \varOb$, $\bL_{\varOb}\subseteq X$. \qedhere
  \end{itemize}
\end{proof}

\noindent
On peut désormais déduire le cas des cardinaux successeurs.

\begin{lemma}
  Soit $\varCk$ un cardinal infini successeur et $\varCk = \SsC{\varCCl}$. On a
  l'inclusion $\bH_\varCk \subseteq \bL_\varCk$.
\end{lemma}

\begin{proof}
  Soit $x \in \bH_\varCk$. On pose d'abord $y = \trcl(\{x\})$, dont on sait que
  $x \in y$ et $\Card(y) \leq \varCCl < \SsC{\varCCl}$. On définit $\rk(y)$
  comme étant le rang minimal tel que $y\subseteq\bL_{\rk(y)}$. On peut trouver
  un cardinal régulier indénombrable $\varCk_y$ tel que
  $\varCk > \varCk_y > \rk(y)$ dans $\bV$, donc dans $\bL$ puisqu'on suppose
  $\VeqL$. On peut vérifier qu'alors,
  $\bL_{\varCk_y}\satisf \ZF - \AxPow + (\VeqL)$. On utilise alors le
  théorème de Löwenheim-Skolem pour trouver un modèle $\varMM$ vérifiant les
  mêmes énoncés de $\SignP{\ZF}$ que $\bL_{\varCk_y}$, contenant $y$ et tel que
  $\Card(\varMM) \leq \varCCl$ (ce qui est possible car $|y|\leq \varCCl$).

  En utilisant alors le \cref{lem.most}, on trouve un ensemble $\varMN$ et un
  isomorphisme $\varCol : (\varMM,\in)\cong (\varMN,\in)$. Comme $y\in \varMM$
  est transitif, il est un point fixe par $\varCol$, donc pour tout
  $z\in y, \varCol(z) = z$. En particulier, $\varCol(x) = x$ et $x\in \varMN$.

  Par le lemme précédent, on sait que $\varMN = \bL_{\varOb}$ pour
  $\varOb = \varMN\cap\Ord \leq \varCk_y < \varCk$, donc (comme $x \in \varMN$)
  $x \in \bL_{\varCk}$.
\end{proof}

\noindent
Le reste de la preuve est une simple vérification.

\begin{lemma}
  Soit $\varCk$ un cardinal infini, on a l'égalité $\bH_\varCk = \bL_\varCk$.
\end{lemma}

\begin{proof}
  Comme on l'a dit, le cas des cardinaux limite est direct par continuité de
  $\bH$ et $\bL$, et le lemme précédent nous donne une inclusion. Il nous reste
  donc à prouver que $\bL_\varCk\subseteq\bH_\varCk$ pour $\varCK$ successeur.

  Si $x\in \bL_\varCk$, alors on trouve $\varOa < \varCk$ tel que $x\in \bLa$,
  mais alors $\trcl(x)\subseteq \bLa$ et, de plus,
  $\Card(\bLa) = \Card(\varOa) < \varCk$.
  Donc $\trcl(x)\in\bH_\varCk$.
\end{proof}

\noindent
On en déduit $\HCG$ dans $\VeqL$.

\begin{theorem}
  $\HCG$ est vérifiée.
\end{theorem}

\begin{proof}
  Soit un cardinal infini $\varCCl$, alors
  $\powerset(\varCCl) \subseteq \bH_{\SsC{\varCCl}} = \bL_{\SsC{\varCCl}}$,
  donc $|\powerset(\varCCl)| \leq \SsC{\lambda}$.
\end{proof}

Nous concluons donc ce chapitre sur ce résultat important~:
\[\Coher(\ZF) \implies \Coher(\ZFC + \HCG)\]

$\HCG$ ne peut donc pas être prouvée faux dans $\ZFC$. Pour montrer que ce
résultat est indépendant de $\ZFC$, il nous reste alors à montrer que
$\Coher(\ZFC)\implies\Coher(\ZFC+\lnot\HC)$, ce qui nécessite la technique du
forcing, introduite par Cohen dans son article
\cite{doi:10.1073/pnas.50.6.1143}.
